% Options for packages loaded elsewhere
\PassOptionsToPackage{unicode}{hyperref}
\PassOptionsToPackage{hyphens}{url}
\documentclass[
  11pt,
]{article}
\usepackage{xcolor}
\usepackage[left=3cm,right=5cm,top=2cm,bottom=2cm]{geometry}
\usepackage{amsmath,amssymb}
\setcounter{secnumdepth}{5}
\usepackage{iftex}
\ifPDFTeX
  \usepackage[T1]{fontenc}
  \usepackage[utf8]{inputenc}
  \usepackage{textcomp} % provide euro and other symbols
\else % if luatex or xetex
  \usepackage{unicode-math} % this also loads fontspec
  \defaultfontfeatures{Scale=MatchLowercase}
  \defaultfontfeatures[\rmfamily]{Ligatures=TeX,Scale=1}
\fi
\usepackage{lmodern}
\ifPDFTeX\else
  % xetex/luatex font selection
\fi
% Use upquote if available, for straight quotes in verbatim environments
\IfFileExists{upquote.sty}{\usepackage{upquote}}{}
\IfFileExists{microtype.sty}{% use microtype if available
  \usepackage[]{microtype}
  \UseMicrotypeSet[protrusion]{basicmath} % disable protrusion for tt fonts
}{}
\makeatletter
\@ifundefined{KOMAClassName}{% if non-KOMA class
  \IfFileExists{parskip.sty}{%
    \usepackage{parskip}
  }{% else
    \setlength{\parindent}{0pt}
    \setlength{\parskip}{6pt plus 2pt minus 1pt}}
}{% if KOMA class
  \KOMAoptions{parskip=half}}
\makeatother
\usepackage{graphicx}
\makeatletter
\newsavebox\pandoc@box
\newcommand*\pandocbounded[1]{% scales image to fit in text height/width
  \sbox\pandoc@box{#1}%
  \Gscale@div\@tempa{\textheight}{\dimexpr\ht\pandoc@box+\dp\pandoc@box\relax}%
  \Gscale@div\@tempb{\linewidth}{\wd\pandoc@box}%
  \ifdim\@tempb\p@<\@tempa\p@\let\@tempa\@tempb\fi% select the smaller of both
  \ifdim\@tempa\p@<\p@\scalebox{\@tempa}{\usebox\pandoc@box}%
  \else\usebox{\pandoc@box}%
  \fi%
}
% Set default figure placement to htbp
\def\fps@figure{htbp}
\makeatother
% definitions for citeproc citations
\NewDocumentCommand\citeproctext{}{}
\NewDocumentCommand\citeproc{mm}{%
  \begingroup\def\citeproctext{#2}\cite{#1}\endgroup}
\makeatletter
 % allow citations to break across lines
 \let\@cite@ofmt\@firstofone
 % avoid brackets around text for \cite:
 \def\@biblabel#1{}
 \def\@cite#1#2{{#1\if@tempswa , #2\fi}}
\makeatother
\newlength{\cslhangindent}
\setlength{\cslhangindent}{1.5em}
\newlength{\csllabelwidth}
\setlength{\csllabelwidth}{3em}
\newenvironment{CSLReferences}[2] % #1 hanging-indent, #2 entry-spacing
 {\begin{list}{}{%
  \setlength{\itemindent}{0pt}
  \setlength{\leftmargin}{0pt}
  \setlength{\parsep}{0pt}
  % turn on hanging indent if param 1 is 1
  \ifodd #1
   \setlength{\leftmargin}{\cslhangindent}
   \setlength{\itemindent}{-1\cslhangindent}
  \fi
  % set entry spacing
  \setlength{\itemsep}{#2\baselineskip}}}
 {\end{list}}
\usepackage{calc}
\newcommand{\CSLBlock}[1]{\hfill\break\parbox[t]{\linewidth}{\strut\ignorespaces#1\strut}}
\newcommand{\CSLLeftMargin}[1]{\parbox[t]{\csllabelwidth}{\strut#1\strut}}
\newcommand{\CSLRightInline}[1]{\parbox[t]{\linewidth - \csllabelwidth}{\strut#1\strut}}
\newcommand{\CSLIndent}[1]{\hspace{\cslhangindent}#1}
\setlength{\emergencystretch}{3em} % prevent overfull lines
\providecommand{\tightlist}{%
  \setlength{\itemsep}{0pt}\setlength{\parskip}{0pt}}
% Three-part paragraph scaffolding for manuscript revision.
%
% bullets  -- boiled-down topic sentence + bullet points
%             small, sans-serif, dark gray; paragraph number [N] prepended
% rgt      -- author's rewrite (initially a placeholder)
%             small, italic, dark blue
% orig     -- original Claude-authored draft
%             normal size, light gray
%
% Presentation order per paragraph: bullets -> rgt -> orig
% Strip with: strip-claudecode -i report.Rmd

\usepackage{xcolor}

\newcounter{pgraph}

\newenvironment{bullets}{%
  \stepcounter{pgraph}%
  \begin{quote}\small\sffamily\color[gray]{0.30}%
  {\normalfont\scriptsize\color[gray]{0.58}[\thepgraph]\enspace}}{%
  \end{quote}}

\newenvironment{rgt}{%
  \begin{quote}\small\itshape\color[RGB]{20,60,160}}{%
  \end{quote}}

\newenvironment{orig}{%
  \begin{quote}\normalsize\color[gray]{0.55}}{%
  \end{quote}}

% backward compat alias
\newenvironment{claudecode}{\begin{orig}}{\end{orig}}
\usepackage{amsmath}
\usepackage{amssymb}
\usepackage{booktabs}
\usepackage{lineno}
\linenumbers
\usepackage{setspace}
\setstretch{1.1}
\usepackage{bookmark}
\IfFileExists{xurl.sty}{\usepackage{xurl}}{} % add URL line breaks if available
\urlstyle{same}
% fallback for those not using the hyperref driver hyperxmp:
\makeatletter
\@ifundefined{xmpquote}{\newcommand{\xmpquote}[1]{#1}}{}
\makeatother
\hypersetup{
  pdftitle={Two Architectures for Simulating Biomarker-Treatment Interactions: Implications for Statistical Power Under Carryover},
  pdfauthor={pmsimstats team},
  hidelinks,
  pdfcreator={LaTeX via pandoc}}

\title{Two Architectures for Simulating Biomarker-Treatment
Interactions: Implications for Statistical Power Under Carryover}
\author{pmsimstats team}
\date{2026-07-31}

\begin{document}
\maketitle

{
\setcounter{tocdepth}{2}
\tableofcontents
}
\section{Abstract}

\textbf{Background.} Simulation studies of predictive
biomarker-treatment interactions in N-of-1 clinical trials require a
data-generating process (DGP) that specifies how the biomarker moderates
treatment response. Two fundamentally different DGP architectures are
used in practice, and the choice between them determines the extent to
which carryover effects reduce statistical power to detect the
interaction. The published simulation literature does not distinguish
them.

\textbf{Methods.} We formalize two architectures: direct mean moderation
(Architecture\textasciitilde A), under which the biomarker scales the
treatment effect in the population mean structure; and differential
correlation (Architecture\textasciitilde B, the multivariate normal
approach), under which the interaction emerges from
treatment-state-dependent correlation between biomarker and response in
the joint distribution. Power consequences under both architectures are
quantified by Monte Carlo simulation on an aggregated N-of-1 design
framework with continuous-time AR(1) residual correlation in the
analysis model. The biomarker, treatment phase, and carryover decay are
parameterized to the prazosin/PTSD reference application.

\textbf{Results.} Under Architecture\textasciitilde A, power is largely
preserved under carryover because the proportional relationship between
biomarker and drug response is maintained at every timepoint. Under
Architecture\textasciitilde B, carryover erodes the differential
correlation signal and reduces power. At the prazosin-calibrated
reference cell (\(N = 35\) per randomization path, i.e.~effective
analysis \(N = 70\) for the crossover (CO) and
open-label-with-blinded-discontinuation (OL+BDC) designs and 140 for
Hybrid after path concatenation; \(c_{bm} = 0.45\), OL+BDC design)
across \(t_{1/2} \in \{0, 0.5, 1.0\}\) weeks,
Architecture\textasciitilde B power drops from \(0.777\) to \(0.539\) (a
30.6\% relative loss; \(z = 7.64\), \(p < 10^{-13}\) on the B-versus-A
contrast), while Architecture\textasciitilde A drops only from \(0.751\)
to \(0.730\) (a 2.8\% relative loss). The B-versus-A absolute power-loss
gap of 21.7 percentage points at OL+BDC, \(t_{1/2} = 1.0\) is
approximately 11 times the matched Architecture\textasciitilde A loss;
the Hybrid design produces a smaller but still significant gap
(\(z = 3.96\), \(p < 10^{-4}\)); the traditional crossover design
produces no detectable gap because its within-subject AR(1) structure
dominates the carryover-mediated correlation channel that
Architecture\textasciitilde B uses. The 36-cell production factorial (2
architectures \(\times\) 3 designs \(\times\) 3 \(t_{1/2}\) \(\times\) 2
\(c_{bm}\) levels) was run at 1\{,\}000 replicates per cell with 100\%
convergence; the Monte Carlo standard error (MCSE) on cell power is
0.013-0.016. Type\textasciitilde I error sits in \([0.010, 0.074]\),
conservative relative to the nominal level and with no
architecture-specific inflation. The conservatism reflects the
\texttt{corCAR1} working-covariance miscalibration characterized in a
companion calibration
study\textsuperscript{\citeproc{ref-pmsimstats-paper10}{1}}; being
common to both architectures, it does not affect the architecture
contrast.

\textbf{Conclusions.} The two architectures encode different biological
mechanisms, heterogeneity at the mean structure versus heterogeneity at
the variance structure, and the choice between them is a substantive
scientific question rather than a parameterization convenience. We
recommend explicit declaration of the DGP architecture in any
biomarker-validation simulation study and provide guidance for selecting
the architecture appropriate to a given clinical context.

\bigskip\noindent

\rule{\textwidth}{0.4pt}\bigskip

\section{1. Introduction}

\begin{bullets}

\textbf{The central question is biomarker-treatment interaction;
simulation studies need a DGP that encodes it.}

\begin{itemize}\setlength{\itemsep}{2pt}\setlength{\parskip}{0pt}
\item Does a baseline biomarker predict differential treatment
  response, i.e., does it moderate the treatment effect?
\item Closed-form power is available only in simplified cases
  [@pmsimstats-paper08]; otherwise simulation is the principal power
  evaluation tool for proposed trial designs.
\item These simulations require a DGP that explicitly encodes
  the biomarker-treatment interaction mechanism.
\end{itemize}

\end{bullets}

\begin{rgt}
rgt to complete.

\end{rgt}

\begin{orig}
A central question in precision medicine trial design is whether a
baseline biomarker predicts differential treatment response, that is,
whether the treatment effect is moderated by the biomarker value. For
simplified special cases, such as a balanced strict repeated-measures
ANOVA under sphericity, the power of the interaction test admits a
closed-form expression, derived together with its equivalence to a
two-sample \(t\)-test in a companion
paper\textsuperscript{\citeproc{ref-pmsimstats-paper08}{2}}. For the
continuous-time mixed model and the N-of-1 designs considered here,
however, no such closed form is available, so simulation studies are the
principal tool for evaluating the statistical power of proposed trial
designs to detect such interactions, and these simulations necessarily
require a data-generating process (DGP) that encodes the
biomarker-treatment interaction mechanism.

\end{orig}

\begin{bullets}

\textbf{Every simulation framework must choose how to generate
treatment effect heterogeneity, and that choice has
consequences.}

\begin{itemize}\setlength{\itemsep}{2pt}\setlength{\parskip}{0pt}
\item Heterogeneity of treatment effects (HTE) literature (2005, 2020)
  distinguishes qualitative interactions (direction reverses)
  from quantitative ones (magnitude varies).
\item Any simulation framework must decide how to encode this
  heterogeneity in the DGP.
\item That decision is substantive, not a mere parameterization
  detail.
\end{itemize}

\end{bullets}

\begin{rgt}
rgt to complete.

\end{rgt}

\begin{orig}
The statistical literature on treatment effect heterogeneity, developed
notably by Rothwell\textsuperscript{\citeproc{ref-rothwell2005}{3}} and
synthesized more recently by Kent et
al.\textsuperscript{\citeproc{ref-kent2020path}{4}}, distinguishes
between \textit{qualitative} interactions, in which the direction of the
treatment effect reverses across biomarker strata, and
\textit{quantitative} interactions, in which the magnitude varies but
the direction remains consistent. As we shall argue, any simulation
framework intended to evaluate power for biomarker-moderated effects
must choose how to generate this heterogeneity, and that choice is not
merely a technical convenience.

\end{orig}

\begin{bullets}

\textbf{Architecture A (direct mean moderation) is the
near-universal standard; the Hendrickson (2020) multivariate
normal (MVN) approach is apparently unique.}

\begin{itemize}\setlength{\itemsep}{2pt}\setlength{\parskip}{0pt}
\item Every major simulation framework uses direct mean
  moderation: Simon (2010) for stratified designs; Freidlin and
  Korn (2014) for adaptive enrichment; Renfro and Sargent
  (2017) for basket/umbrella trials; all N-of-1 frameworks.
\item Wang and Schork (2019) and Schork (2022) also use
  Architecture A despite sharing authorship with Hendrickson
  et al. (2020), suggesting the MVN parameterization is a
  paper-specific choice, not a lab convention.
\item To the best of our knowledge, no other clinical trial
  simulation framework uses differential correlation (MVN) to
  encode the biomarker-treatment interaction.
\end{itemize}

\end{bullets}

\begin{rgt}
rgt to complete.

\end{rgt}

\begin{orig}
The dominant approach in the clinical trial simulation literature is
direct mean moderation, in which the biomarker enters the outcome model
as an interaction term that explicitly scales the treatment effect.
Simon\textsuperscript{\citeproc{ref-simon2010}{5}} employs this approach
throughout the simulation frameworks he develops for
biomarker-stratified designs; Freidlin and
Korn\textsuperscript{\citeproc{ref-freidlin2014}{6}} adopt it for
adaptive enrichment trials; Renfro and
Sargent\textsuperscript{\citeproc{ref-renfro2017}{7}} use it throughout
their treatment of basket and umbrella trials; and it underlies the
platform trial literature without exception. Recent simulation work on
biomarker-treatment interaction estimation in randomized trials with
prognostic covariates, specifically the study by Haller and
Ulm\textsuperscript{\citeproc{ref-haller2018}{8}}, likewise adopts the
standard regression specification with explicit interaction terms. Among
N-of-1 simulation studies, Zucker et
al.\textsuperscript{\citeproc{ref-zucker1997}{9}}, Araujo et
al.\textsuperscript{\citeproc{ref-araujo2016}{10}}, and Duan et al.
(2013) similarly use hierarchical random-effects models in which the
biomarker predicts the individual treatment effect through the mean
structure. Notably, Wang and
Schork\textsuperscript{\citeproc{ref-wang2019}{11}} and
Schork\textsuperscript{\citeproc{ref-schork2022}{12}}, whose
methodological work addresses power and design for N-of-1 trials with
serial correlation and carryover directly, also specify the treatment
effect through a linear mean structure, which suggests that the
multivariate-normal differential-correlation parameterization of
Hendrickson et al.\textsuperscript{\citeproc{ref-hendrickson2020}{13}}
is not a convention of any particular research program but rather a
paper-specific design choice. To the best of our knowledge, no other
clinical trial simulation framework generates the biomarker-treatment
interaction through differential correlation in a multivariate normal
joint distribution rather than through direct mean moderation.

\end{orig}

\begin{bullets}

\textbf{Architecture B is principled, not merely idiosyncratic: it
is the reduced-form second-moment encoding of a latent
responder-class mechanism.}

\begin{itemize}\setlength{\itemsep}{2pt}\setlength{\parskip}{0pt}
\item If the biomarker is a proxy for an unobserved responder
  subtype, the faithful DGP is a finite mixture; a single MVN
  reproduces its first two moments.
\item Architecture B is thus one point on a continuum from
  continuous-heterogeneity to explicit latent-class formulations.
\item When to prefer it, and whether the approximation is adequate,
  is taken up in Section 4.2 and formally in a companion paper
  [@pmsimstats-paper03].
\end{itemize}

\end{bullets}

\begin{rgt}
rgt to complete.

\end{rgt}

\begin{orig}
Although the differential-correlation parameterization is unusual in the
simulation literature, it is not arbitrary. It is the reduced-form,
second-moment representation of a latent responder-class mechanism: if
the biomarker does not itself determine the drug's effect but instead
indexes an unobserved subpopulation that responds (a disease subtype, an
inflammatory state, a genetic background), the biologically faithful DGP
is a finite mixture, and a single multivariate normal reproduces that
mixture's first two moments through a treatment-state-dependent
biomarker-response correlation. On this reading
Architecture\textasciitilde B is not a competitor to mean moderation on
equal footing but the covariance-only end of a continuum that runs
through mean-and-covariance and location-scale specifications to an
explicit latent-class model. When this reduced form is preferable to
direct mean moderation, and when the approximation breaks down, is taken
up in Section\textasciitilde4.2 and is the subject of a companion
paper\textsuperscript{\citeproc{ref-pmsimstats-paper03}{14}} that
develops this biomarker-moderation spectrum and evaluates class-aware
analyses at trial-relevant sample sizes.

In plain language, and without the latent-variable vocabulary: suppose
patients come in two kinds, those who respond to the drug and those who
do not, but we cannot see which kind any given patient is. The biomarker
is only an imperfect clue to that hidden status, so two patients with
the same biomarker value can still respond differently.
Architecture\textasciitilde A cannot capture this, because it assumes
the biomarker fixes the size of each patient's drug effect exactly.
Architecture\textasciitilde B captures it by letting the biomarker and
the response track each other only while the drug is acting: a higher
biomarker tends to go with a larger response, as a statistical tendency
rather than a fixed rule. When the drug washes out that tendency fades,
which is the reason carryover costs more power under
Architecture\textasciitilde B than under Architecture\textasciitilde A.

\end{orig}

\begin{bullets}

\textbf{The comparison originates not in an error but in
recognizing that Hendrickson's Architecture B is appropriate for
its latent-subtype biomarker, whereas the N-of-1 simulation
literature defaults to Architecture A.}

\begin{itemize}\setlength{\itemsep}{2pt}\setlength{\parskip}{0pt}
\item Hendrickson et al. (2020) use Architecture B, the correct
  encoding for a biomarker hypothesized to proxy a latent responder
  subtype (resting systolic blood pressure for a noradrenergic PTSD
  subtype).
\item Most N-of-1 and biomarker-trial simulation work instead
  encodes the interaction through Architecture A.
\item The divergence became concrete when two implementations of
  the same pmsimstats pipeline, one using B and one A, reached
  different conclusions about carryover.
\end{itemize}

\end{bullets}

\begin{rgt}
rgt to complete.

\end{rgt}

\begin{orig}
Architecture\textasciitilde B is the appropriate data-generating choice
for the trial that motivated Hendrickson et
al.\textsuperscript{\citeproc{ref-hendrickson2020}{13}}. The predictive
biomarker there is resting systolic blood
pressure\textsuperscript{\citeproc{ref-raskind2013}{15}}, which is
hypothesized to index a latent noradrenergic subtype of PTSD
preferentially responsive to alpha-1 adrenergic
blockade\textsuperscript{\citeproc{ref-raskind2018}{16}} rather than to
determine prazosin exposure. On that hypothesis the biomarker is a proxy
for an unobserved responder class, exactly the latent responder-class
mechanism for which the differential-correlation encoding is the reduced
form (Section\textasciitilde4.3). What brought the architectural
distinction into focus was the realization that this choice is not
shared by the wider N-of-1 and biomarker-trial simulation literature,
which, as reviewed above, encodes the interaction almost exclusively
through direct mean moderation (Architecture\textasciitilde A). The
mismatch became concrete during development of the pmsimstats framework:
the original Hendrickson-derived code generated data under
Architecture\textasciitilde B, while a parallel reimplementation adopted
Architecture\textasciitilde A, and the two produced qualitatively
different conclusions about how carryover erodes power. The distinction
is therefore not a coding artifact to be corrected but a substantive
modeling choice. The same biomarker-moderation question, encoded in the
first versus the second moment of the outcome distribution, yields
materially different power under carryover, and the encoding that is
appropriate depends on the biology of the biomarker.

\end{orig}

\begin{bullets}

\textbf{This paper is the first systematic comparison of the
two DGP architectures; we formalize, simulate, and give
guidance.}

\begin{itemize}\setlength{\itemsep}{2pt}\setlength{\parskip}{0pt}
\item We formalize the distinction between Architecture A
  (mean moderation) and Architecture B (differential
  correlation).
\item We demonstrate power consequences through Monte Carlo
  simulation.
\item We offer guidance for investigators on how to choose
  the appropriate architecture.
\end{itemize}

\end{bullets}

\begin{rgt}
rgt to complete.

\end{rgt}

\begin{orig}
We present here what we believe to be the first systematic comparison of
these two DGP architectures for biomarker-treatment interaction power
estimation under carryover. We shall formalize the distinction between
them, demonstrate its consequences through Monte Carlo simulation, and
offer guidance for investigators designing simulation studies for
predictive biomarker-moderated treatment effects.

\end{orig}

\textbf{Related documents in this series.} The mathematical derivation,
code-level implementation details, four implementation attempts (three
failed and one correct), and the empirical evaluation of the original
Hendrickson\textsuperscript{\citeproc{ref-hendrickson2020}{13}}
parameter grid under Architecture\textasciitilde A are documented in
\texttt{19-mean-moderation-implementation-notes.tex} (with a brief
markdown summary in
\texttt{19-mean-moderation-implementation-notes.md}). A condensed
empirical-results summary suitable for collaborator presentation is
provided in \texttt{21-architecture-comparison-summary.md}.

\bigskip\noindent

\rule{\textwidth}{0.4pt}\bigskip

\section{2. Two DGP Architectures}

\textbf{Notation.} Throughout this paper, \(B_i\) denotes the biomarker
value of participant \(i\), \(D_{it}\) the binary on-drug indicator at
timepoint \(t\) (1 on drug, 0 off drug), and \(BR_{it}\) the biological
response component of the outcome. These three symbols describe the
data-generating-process layer. The analysis-model layer uses a related
but distinct continuous drug exposure variable, \texttt{Dbc}, defined in
Section 3.1. Code-style identifiers in typewriter font (\texttt{bm},
\texttt{Dbc}, \texttt{BR}) refer to specific variables in the pmsimstats
codebase or to R formula syntax; mathematical symbols are used
otherwise.

\subsection{2.1 Architecture A: Direct Mean Moderation}

In this approach the biomarker enters the mean structure of the outcome
directly. The treatment effect for participant \(i\) at timepoint \(t\)
is:

\[Y_{it} = \mu_0 + \beta_1 \cdot D_{it} \cdot (1 + \beta_{bm}
\cdot B_i) + \epsilon_{it}\]

where \(D_{it}\) is the binary on-drug indicator (or a continuous
measure of drug exposure), \(B_i\) is the participant's biomarker value
(typically centered), and \(\beta_{bm}\) is the biomarker moderation
parameter. The interaction is explicit in the mean: participants with
higher biomarker values have proportionally larger treatment effects.

Key properties:

\begin{itemize}
\item The interaction is \textbf{deterministic} given the biomarker value
  and treatment status.
\item The biomarker moderation $\beta_{bm}$ is a \textbf{population-level
  parameter} that applies uniformly to all participants.
\item The signal for the interaction is in the \textbf{first moment}
  (conditional mean) of the outcome distribution.
\item Adding noise, random effects, or carryover to $D_{it}$ scales
  the entire treatment effect (including its biomarker-dependent
  component) proportionally. The \textbf{ratio} of drug effect between
  high-biomarker and low-biomarker participants is preserved.
\end{itemize}

\begin{bullets}

\textbf{The multiplicative and additive forms of Architecture A
are interchangeable for power calculations.}

\begin{itemize}\setlength{\itemsep}{2pt}\setlength{\parskip}{0pt}
\item The multiplicative form maps directly to the standard
  additive regression, the effect-modeling equation of
  Kent et al. (2020).
\item $\beta_3 = \beta_1 \cdot \beta_{bm}$; the additive form
  additionally allows a prognostic main effect.
\item Both forms fall within Architecture A and are
  interchangeable for power calculations on the interaction
  coefficient.
\end{itemize}

\end{bullets}

\begin{rgt}
rgt to complete.

\end{rgt}

\begin{orig}
We note that the multiplicative form above maps directly onto the
standard additive regression form for predictive biomarker effects used
in the trial methodology literature, specifically the effect modeling
equation of Kent et al.\textsuperscript{\citeproc{ref-kent2020path}{4}}:
\[Y_{it} = \beta_0 + \beta_1 D_{it} + \beta_2 B_i + \beta_3
D_{it} B_i + \epsilon_{it}\] The interaction coefficient \(\beta_3\) in
the additive form corresponds to the product
\(\beta_1 \cdot \beta_{bm}\) in the multiplicative form. The additive
form additionally accommodates a prognostic main effect \(\beta_2\) for
the biomarker, which the multiplicative form sets to zero by
construction. For the purpose of power calculations on the interaction
coefficient, the two forms are interchangeable up to this
prognostic-effect distinction, and both fall within Architecture A.

\end{orig}

This architecture is used in the pedagogical simulation
(\texttt{implementations/simple/simulation.R}) and in the exploratory
carryover analyses (\texttt{simulation\_carryover\_spectrum.R}) of the
pmsimstats project.

\subsection{2.2 Architecture B: Differential Correlation (MVN)}

In this approach the biomarker and the biological response (BR)
component are jointly drawn from a multivariate normal distribution with
treatment-state-dependent correlation:

\[\begin{pmatrix} B_i \\ BR_{it} \end{pmatrix} \sim
\text{MVN}\left(\begin{pmatrix} \mu_B \\ \mu_{BR}(t, D_{it})
\end{pmatrix}, \Sigma(D_{it})\right)\]

where the covariance matrix \(\Sigma\) has the following
treatment-state-dependent BM-BR correlation:

\[\text{Cor}(B_i, BR_{it}) = \begin{cases}
c_{bm} & \text{if } D_{it} = 1 \text{ (on drug)} \\
c_{bm} \cdot e^{-\lambda \cdot t_{sd}} &
\text{if } D_{it} = 0 \text{ (off drug)}
\end{cases}\]

where \(t_{sd}\) is time since drug discontinuation. The
exponential-decay form is the default; the limiting cases
\(\lambda = 0\) (no decay; correlation persists indefinitely during the
off-drug phase) and \(\lambda \to \infty\) (immediate decay; correlation
drops to zero the instant the drug is discontinued) recover the two
no-carryover model variants.

\begin{bullets}

\textbf{In Architecture B the interaction is in the covariance
structure, not the mean; it emerges from the conditional
distribution given the biomarker value.}

\begin{itemize}\setlength{\itemsep}{2pt}\setlength{\parskip}{0pt}
\item The population mean of $BR_{it}$ is the same for all
  participants; no mean-structure interaction exists.
\item Given biomarker value $b$, the conditional expectation
  of BR scales with $c_{bm}$, creating an interaction-like
  structure.
\item Correlation is $c_{bm}$ on drug and decays as
  $c_{bm} e^{-\lambda t_{sd}}$ off drug; moderation attenuates
  during the off-drug period.
\end{itemize}

\end{bullets}

\begin{rgt}
rgt to complete.

\end{rgt}

\begin{orig}
The interaction is not in the mean structure: the population mean of
\(BR_{it}\) is the same for all participants with the same treatment
history. Instead it emerges from the \textbf{conditional distribution}:
given a participant's biomarker value \(B_i = b\), the conditional
expectation of their biological response is:

\[E[BR_{it} | B_i = b, D_{it} = 1] = \mu_{BR}(t) + c_{bm}
\cdot \frac{\sigma_{BR}}{\sigma_B} \cdot (b - \mu_B)\]

This conditional expectation has an interaction-like structure: the
biomarker value \(b\) modulates the expected BR. The effective
correlation is \(c_{bm}\) on drug and \(c_{bm} \cdot
e^{-\lambda t_{sd}}\) during the off-drug period; the moderation
therefore persists with attenuated strength during off-drug periods and
vanishes only in the limit of complete decay.

\end{orig}

Key properties:

\begin{itemize}
\item The interaction is \textbf{probabilistic}: it manifests as a
  tendency, not a deterministic scaling.
\item The biomarker moderation $c_{bm}$ is a \textbf{correlation
  parameter} that governs the strength of the association.
\item The signal for the interaction is in the \textbf{second moment}
  (covariance structure) of the joint distribution.
\item Carryover effects in the DGP inflate the BR means during
  off-drug periods, making them resemble on-drug means. This
  reduces the \textbf{differential} in the correlation structure
  between treatment states, weakening the signal that the
  analysis model detects.
\end{itemize}

This architecture is used in the Hendrickson et
al.\textsuperscript{\citeproc{ref-hendrickson2020}{13}} publication code
and the audited \texttt{R/} package in pmsimstats-ng.

A third possibility is that both channels operate at once: the biomarker
could scale the mean drug effect (Architecture\textasciitilde A) while
simultaneously altering the biomarker-response correlation during active
treatment (Architecture\textasciitilde B). Such a combined architecture,
carrying independent weights on the two channels, recovers A and B as
its single-channel limits and is the natural framework when the
operative mechanism is uncertain. Because it introduces a second free
parameter and its power behavior warrants separate treatment, we develop
it in a companion
paper\textsuperscript{\citeproc{ref-pmsimstats-paper11}{17}} and
restrict the present paper to the A-versus-B contrast.

\subsection{2.3 Mathematical Comparison}

\begin{bullets}

\textbf{Under Architecture A, carryover shrinks the signal
amplitude but preserves the ratio between biomarker groups,
so identifiability is maintained.}

\begin{itemize}\setlength{\itemsep}{2pt}\setlength{\parskip}{0pt}
\item The expected outcome difference between two participants
  scales linearly with $D$.
\item Under carryover, $D$ is replaced by a decayed value and
  the absolute magnitude shrinks.
\item But the ratio of drug effect between high- and
  low-biomarker participants is preserved at every exposure
  level; the signal contracts but remains identifiable.
\end{itemize}

\end{bullets}

\begin{rgt}
rgt to complete.

\end{rgt}

\begin{orig}
We begin the comparison by considering the expected outcome difference
between two participants with biomarker values \(b_1\) and \(b_2\) at
the same drug exposure level \(D\) under Architecture A:

\[E[Y \mid b_1, D] - E[Y \mid b_2, D] = \beta_1 \cdot \beta_{bm}
\cdot (b_1 - b_2) \cdot D\]

The absolute magnitude of this difference scales linearly with \(D\) and
therefore shrinks during off-drug periods when \(D\) is replaced by an
exposure-decayed value \(D \cdot \text{carryover}\). We note, however,
that the \textbf{ratio} of the drug effect between high-biomarker and
low-biomarker participants is preserved at every level of exposure. The
biomarker-treatment signal contracts in amplitude but does not lose its
identifiability under carryover.

\end{orig}

\begin{bullets}

\textbf{Under Architecture B, the biomarker-dependent component
of the off-drug response decays exponentially, eventually
eliminating the differential signal.}

\begin{itemize}\setlength{\itemsep}{2pt}\setlength{\parskip}{0pt}
\item Off-drug, the conditional expectation includes a
  biomarker-dependent term scaled by $c_{bm} e^{-\lambda
  t_{sd}}$.
\item As $t_{sd}$ grows, that term shrinks toward zero.
\item In the limit, all participants converge to the same
  off-drug mean regardless of biomarker value; the interaction
  signal vanishes.
\end{itemize}

\end{bullets}

\begin{rgt}
rgt to complete.

\end{rgt}

\begin{orig}
Under Architecture B, the analogous quantity is:

\[E[BR | b_1, D=1] - E[BR | b_2, D=1] = c_{bm} \cdot
\frac{\sigma_{BR}}{\sigma_B} \cdot (b_1 - b_2)\]

During off-drug periods with carryover:

\[E[BR | b, D=0, t_{sd}] = \mu_{BR,\text{off}}(t_{sd}) +
c_{bm} \cdot e^{-\lambda t_{sd}} \cdot
\frac{\sigma_{BR}}{\sigma_B} \cdot (b - \mu_B)\]

Here the biomarker-dependent component decays exponentially with time
off drug. As \(t_{sd}\) increases, the off-drug conditional expectation
converges to \(\mu_{BR,\text{off}}\) for all biomarker values,
eliminating the differential signal entirely.

\end{orig}

\bigskip\noindent

\rule{\textwidth}{0.4pt}\bigskip

\section{3. Consequences for Statistical Power Under Carryover}

\subsection{3.1 Empirical demonstration}

\begin{bullets}

\textbf{The two DGPs are identical except for how they encode
the BM-BR linkage; any power divergence is attributable solely
to that difference.}

\begin{itemize}\setlength{\itemsep}{2pt}\setlength{\parskip}{0pt}
\item All non-biologic components (time-varying, placebo-belief,
  residual noise) are generated identically in both
  architectures.
\item The only difference: Architecture B encodes BM-BR
  interaction in $\Sigma$; Architecture A adds a post-draw
  mean shift to on-drug BR.
\item The same \texttt{nlme::lme} analysis model with
  \texttt{corCAR1} is used for both, so divergence in results
  is attributable to the DGP alone.
\end{itemize}

\end{bullets}

\begin{rgt}
rgt to complete.

\end{rgt}

\begin{orig}
We ran identical simulation configurations under both architectures
using the pmsimstats-ng framework, matching the DGP parameters as
closely as possible between the two approaches. We note that the
non-biologic response components (time-varying response, placebo-belief
accumulation, and residual noise) are generated from the same
multivariate normal draw with identical means, variances, within-factor
AR(1) autocorrelation, and cross-factor correlations in both
architectures. The two DGPs differ \emph{only} in how the
biomarker-to-biological-response (BM-BR) linkage is encoded:
Architecture\textasciitilde B embeds it in the off-diagonal \(c_{bm}\)
entries of \(\Sigma\), whereas Architecture\textasciitilde A omits those
entries and adds a post-draw shift
\(c_{bm} \cdot B_i^{*} \cdot \sigma_{BR}\) to on-drug BR columns (see
\texttt{19-mean-moderation-implementation-notes.tex} for the scaling
derivation). The divergent carryover behavior reported below is
therefore attributable to the BM-BR parameterization, not to any
difference in how the placebo, time-varying, or residual-noise
components are generated. The analysis model was the same in both cases:
an \texttt{nlme::lme} fit with a \texttt{corCAR1} continuous-time
autocorrelation structure and a single biomarker-by-drug interaction
term, written \texttt{bm:Dbc} in R formula syntax.

\end{orig}

\begin{bullets}

\textbf{\texttt{Dbc} is a continuous drug exposure variable
that captures exponential off-drug decay; the \texttt{bm:Dbc}
interaction tests whether biomarker moderates that
exposure-decayed effect.}

\begin{itemize}\setlength{\itemsep}{2pt}\setlength{\parskip}{0pt}
\item \texttt{Dbc} = 1 on drug; = $\exp(-\lambda t_{sd})$ off
  drug, matching the DGP decay rate.
\item Unlike the binary $D_{it}$, \texttt{Dbc} smoothly
  captures off-drug shrinkage.
\item The test of \texttt{bm:Dbc} is the test of whether the
  biomarker moderates the (decay-adjusted) drug effect.
\end{itemize}

\end{bullets}

\begin{rgt}
rgt to complete.

\end{rgt}

\begin{orig}
The analysis-model drug exposure variable \texttt{Dbc} is the
\emph{continuous} analogue of the binary DGP indicator \(D_{it}\). It
equals \(1\) during on-drug periods and \(\exp(-\lambda t_{sd})\) during
off-drug periods, where \(t_{sd}\) is time since drug discontinuation
and \(\lambda\) is the analysis-model carryover decay rate (typically
derived from the drug's pharmacokinetic half-life as
\(\lambda = \ln 2 / t_{1/2}\), in which case it matches the DGP decay
rate of Section 2.2 by construction). Whereas \(D_{it}\) flips abruptly
between 0 and 1, \texttt{Dbc} captures the smooth shrinkage of residual
drug effect during the off-drug period. The interaction coefficient on
\texttt{bm:Dbc} is therefore the test of whether the biomarker moderates
the exposure-decayed drug effect.

\end{orig}

We follow the ADEMP framework of Morris, White, and Crowther (2019):
Aims, Data-generating mechanisms, Estimands, Methods, Performance
measures.

\textbf{Setup:}

\begin{itemize}
\item Trial designs: CO (traditional crossover), Hybrid (N-of-1
  Hybrid), OL+BDC (open-label with blinded discontinuation)
\item Sample size: $N = 35$ per randomization path (effective
  analysis $N = 70$ for CO and OL+BDC, 140 for Hybrid, after the
  paths are concatenated)
\item Biomarker correlation / moderation: $c_{bm} = 0.45$ /
  $\beta_{bm} = 0.45$
\item DGP carryover half-life: $t_{1/2} \in \{0, 0.5, 1.0\}$ weeks
\item Analysis: matched \texttt{Dbc} with same half-life
\item Replicates: 1{,}000 per cell (MCSE
  on a binomial power estimate near $0.5$ is approximately
  $\sqrt{0.5 \cdot 0.5 / 1{,}000} \approx 0.016$)
\item Implementation: 8-core \texttt{parallel::mclapply}; full
  factorial 2 architectures $\times$ 3 designs $\times$ 3
  $t_{1/2}$ $\times$ 2 $c_{bm}$ levels = 36 cells, 36{,}000
  fits, 100\% convergence
\end{itemize}

\begin{bullets}

\textbf{$N = 35$ was chosen because it keeps power off the
ceiling, making any carryover-induced loss visible; Type I
error is in the expected range.}

\begin{itemize}\setlength{\itemsep}{2pt}\setlength{\parskip}{0pt}
\item $N = 35$ is the smallest prazosin-PTSD-calibrated sample
  size in the literature.
\item Neither architecture saturates power at zero carryover,
  so losses are mechanically observable.
\item Type I error across 18 null cells sat in
  $[0.010, 0.074]$, consistent with the conservative-to-nominal
  range expected for this specification.
\end{itemize}

\end{bullets}

\begin{rgt}
rgt to complete.

\end{rgt}

\begin{orig}
The \(N = 35\) choice is the smallest of the prazosin-PTSD- calibrated
sample sizes used in the published methodological literature; it is a
regime in which neither architecture saturates power at the no-carryover
cells, so any carryover-induced erosion is mechanically visible. Type I
error under the \(c_{bm} = 0\) null cells sat in \([0.010, 0.074]\)
across the 18 null cells. These rates below the nominal 0.05 are not
architecture-specific; they reflect a known conservatism of the standard
\texttt{corCAR1} Wald test, whose model-based standard error is
over-stated by roughly 6 to 10 percent under this working-covariance
specification. A companion calibration
study\textsuperscript{\citeproc{ref-pmsimstats-paper10}{1}} localizes
this miscalibration to the \texttt{corCAR1} residual layer and shows
that a CR2 bias-reduced cluster-robust standard error restores the
realized Type I error to approximately the nominal level, with a
near-uniform power increment that leaves the ranking of analysis
specifications unchanged. Because the conservatism is common to both
architectures, it is common-mode with respect to the B-versus-A contrast
that is the object of this section and does not affect it.

\end{orig}

\textbf{Results under Architecture A (direct mean moderation):}

\begin{table}[htbp]
\centering
\begin{tabular}{lccc}
\toprule
Design & $t_{1/2} = 0$ & $t_{1/2} = 0.5$ & $t_{1/2} = 1.0$ \\
\midrule
CO & 0.739 & 0.739 & 0.725 \\
Hybrid & 0.992 & 0.987 & 0.978 \\
OL+BDC & 0.751 & 0.748 & 0.730 \\
\bottomrule
\end{tabular}
\end{table}

\begin{bullets}

\textbf{Architecture A shows modest power loss under carryover
(1--3\%), below preliminary expectations from earlier exploratory
runs in this project.}

\begin{itemize}\setlength{\itemsep}{2pt}\setlength{\parskip}{0pt}
\item Relative losses: 1.9\% (CO), 1.4\% (Hybrid), 2.8\%
  (OL+BDC).
\item All three fall below the 8--13\% range suggested by earlier
  exploratory runs in this project.
\item Qualitative finding: modest erosion; power essentially
  intact under carryover.
\end{itemize}

\end{bullets}

\begin{rgt}
rgt to complete.

\end{rgt}

\begin{orig}
Power declines modestly under Architecture A: relative loss from
\(t_{1/2} = 0\) to \(t_{1/2} = 1\) is 1.9\% (CO), 1.4\% (Hybrid), and
2.8\% (OL+BDC). We note that all three values fall well below the 8-13\%
relative-loss range suggested by earlier exploratory runs in this
project for direct-mean-moderation DGPs, though they reproduce the
qualitative finding of those runs, namely a modest erosion that leaves
power essentially intact.

\end{orig}

\textbf{Results under Architecture B (differential correlation,
MVN):}

\begin{table}[htbp]
\centering
\begin{tabular}{lccc}
\toprule
Design & $t_{1/2} = 0$ & $t_{1/2} = 0.5$ & $t_{1/2} = 1.0$ \\
\midrule
CO & 0.745 & 0.754 & 0.723 \\
Hybrid & 0.995 & 0.980 & 0.944 \\
OL+BDC & 0.777 & 0.694 & 0.539 \\
\bottomrule
\end{tabular}
\end{table}

\begin{bullets}

\textbf{Architecture B shows sharply larger power loss under
carryover, especially OL+BDC (30.6\%); the B-vs-A gap is
significant for OL+BDC and Hybrid but not CO.}

\begin{itemize}\setlength{\itemsep}{2pt}\setlength{\parskip}{0pt}
\item Relative losses: 3.0\% (CO), 5.1\% (Hybrid), 30.6\%
  (OL+BDC); the OL+BDC loss is $\approx$11$\times$ the
  matched Architecture A loss.
\item B-vs-A gap: $z = 7.64$ ($p < 10^{-13}$) at OL+BDC;
  $z = 3.96$ ($p < 10^{-4}$) at Hybrid.
\item CO gap not significant ($z = 0.29$): the CO design's
  within-subject AR(1) structure dominates the
  correlation-decay channel that Architecture B depends on.
\end{itemize}

\end{bullets}

\begin{rgt}
rgt to complete.

\end{rgt}

\begin{orig}
Power declines more sharply under Architecture B, and most markedly so
in the OL+BDC design: relative loss from \(t_{1/2} = 0\) to
\(t_{1/2} = 1\) is 3.0\% (CO; not significantly different from
Architecture A's CO loss), 5.1\% (Hybrid), and \textbf{30.6\% (OL+BDC)}.
The CO cell shows a small, non-monotone fluctuation
(\(0.745 \to 0.754 \to 0.723\)), which is within one MCSE and is
consistent with the null finding of no significant
architecture-by-carryover interaction in that design. The OL+BDC loss is
approximately 11 times the matched Architecture-A OL+BDC loss (2.8\%).
It falls below the 40-60\% relative-loss range suggested by earlier
exploratory runs in this project for the MVN differential-correlation
DGP, but it confirms the qualitative finding of those runs, namely a
substantial and strongly design-dependent erosion. We compare the
carryover-induced power loss between architectures using a
difference-of-differences \(z\)-statistic. Let
\(\Delta_X = \hat{\pi}_{X,0} - \hat{\pi}_{X,1}\) denote the
within-architecture loss from \(t_{1/2} = 0\) to \(t_{1/2} = 1.0\) for
architecture \(X \in \{A, B\}\), estimated from \(n = 1{,}000\)
independent replicates at each cell. The B-versus-A gap is
\(\Delta_B - \Delta_A\), with standard error

\[
\text{SE} = \sqrt{
  \frac{\hat{\pi}_{B,0}(1-\hat{\pi}_{B,0})}{n}
+ \frac{\hat{\pi}_{B,1}(1-\hat{\pi}_{B,1})}{n}
+ \frac{\hat{\pi}_{A,0}(1-\hat{\pi}_{A,0})}{n}
+ \frac{\hat{\pi}_{A,1}(1-\hat{\pi}_{A,1})}{n}}
\]

since the four proportions come from independent replicate sets. All
tests are two-sided. At OL+BDC this gives \(z = 7.64\)
(\(p < 10^{-13}\)); at Hybrid, \(z = 3.96\) (\(p < 10^{-4}\)); the CO
contrast is not statistically distinguishable from zero (\(z = 0.29\))
because the CO design's within-subject AR(1) structure dominates the
carryover-mediated correlation decay channel that Architecture B uses to
encode the biomarker-treatment interaction.

\end{orig}

\subsection{3.2 Why the architectures diverge}

We turn now to a mechanistic account of why the two architectures
produce such different power losses under carryover.

\begin{bullets}

\textbf{Architecture A is subject to only one carryover
channel (mean blurring); it escapes the second channel
(correlation erosion).}

\begin{itemize}\setlength{\itemsep}{2pt}\setlength{\parskip}{0pt}
\item Carryover compresses \texttt{Dbc} contrast and inflates
  off-drug BR means; both architectures share this channel.
\item Architecture A lacks a second channel because $\beta_{bm}$
  is a fixed population parameter and the noise structure is
  drug-state-independent.
\item Result: reduced variance in the interaction term, but
  signal-to-noise ratio degrades only modestly.
\end{itemize}

\end{bullets}

\begin{rgt}
rgt to complete.

\end{rgt}

\begin{orig}
\textbf{Under Architecture A}, carryover reduces the magnitude of the
drug exposure variable seen by the analysis model (\texttt{Dbc} drops
below 1 during the off-drug period), and the mean BR at off-drug
timepoints is inflated by the residual drug effect exactly as in
Architecture B. Both mechanisms compress the between-state treatment
contrast and are shared by the two architectures. What Architecture A
lacks is the second channel of signal loss: the biomarker moderation
coefficient \(\beta_{bm}\) is a fixed population parameter, the
biomarker enters the mean structure directly rather than through a
correlation, and the noise structure is independent of drug state. The
biomarker-by-drug interaction term therefore has reduced variance
(because \texttt{Dbc} has less contrast between drug states), but the
signal-to-noise ratio degrades only modestly.

\end{orig}

\textbf{Under Architecture B}, carryover operates on two levels
simultaneously:

\begin{enumerate}
\item \textbf{Mean blurring}: The BR means during off-drug periods are
  inflated by residual carryover, making them resemble on-drug
  BR means. This reduces the treatment contrast that the analysis
  model uses to distinguish drug effect from no-drug. This channel
  is shared with Architecture A.

\item \textbf{Correlation erosion}: The BM-BR correlation during off-drug
  periods decays exponentially with $t_{sd}$. This is not just a
  statistical modeling choice; it is a structural consequence
  of the DGP. The correlation between biomarker and BR reflects
  the degree to which the drug effect is active; as the drug
  washes out, the correlation weakens because the drug-mediated
  component of BR variance shrinks relative to the
  drug-independent component. This channel is \emph{unique} to
  Architecture B, because only here does the biomarker signal
  live in the second moment of the joint distribution. Under
  Architecture A the drug-mediated and drug-independent
  components of BR variance move in tandem (both are scaled by
  the same \texttt{Dbc}-weighted drug effect), so their ratio,
  and hence the identifiability of the interaction, is preserved
  even as amplitudes shrink.
\end{enumerate}

\begin{bullets}

\textbf{Under Architecture B the interaction estimate is
attacked from both sides simultaneously; Architecture A faces
only one of the two attacks.}

\begin{itemize}\setlength{\itemsep}{2pt}\setlength{\parskip}{0pt}
\item The \texttt{bm:Dbc} coefficient depends on the covariance
  between $B_i$ and \texttt{Dbc}-weighted outcomes.
\item Under Architecture B: (1) \texttt{Dbc} contrast is
  compressed; (2) the BM-BR correlation generating the
  covariance signal is simultaneously decaying.
\item Architecture A faces only attack (1), not attack (2);
  this single difference accounts for the divergence in power.
\end{itemize}

\end{bullets}

\begin{rgt}
rgt to complete.

\end{rgt}

\begin{orig}
The analysis model's interaction coefficient (the coefficient on
\texttt{bm:Dbc}) depends on the covariance between \(B_i\) and
\texttt{Dbc}-weighted outcomes. Under Architecture B, this covariance is
being attacked from both sides: the treatment contrast in \texttt{Dbc}
is compressed, and the BM-BR correlation that generates the covariance
signal is simultaneously decaying. That Architecture A escapes the
second attack entirely accounts for the divergence in power loss
documented above.

\end{orig}

\subsection{3.3 Robustness check at $N = 70$}

\begin{bullets}

\textbf{An $N = 70$ confirmation run checks whether the
qualitative ordering holds when power is near the ceiling.}

\begin{itemize}\setlength{\itemsep}{2pt}\setlength{\parskip}{0pt}
\item 800 replicates per cell, 12 cells (2 architectures
  $\times$ 2 $c_{bm}$ $\times$ 3 $t_{1/2}$), OL+BDC design.
\item 100\% convergence across 9{,}600 fits.
\item Purpose: verify the Section 3.1 ordering survives
  saturation at the larger $N$.
\end{itemize}

\end{bullets}

\begin{rgt}
rgt to complete.

\end{rgt}

\begin{orig}
The Section\textasciitilde3.1 production results are at \(N = 35\),
where neither architecture saturates power and the carryover-induced
loss is mechanically visible. We conducted a confirmation run at
\(N = 70\) to check whether the qualitative ordering survives the
saturation that comes with the larger sample size. An
800-replicate-per-cell Monte Carlo run was executed at \(N = 70\) across
12 cells (architecture \(\in \{\text{MVN},
\text{mean moderation}\} \times c_{bm} \in \{0, 0.45\} \times
t_{1/2} \in \{0, 0.5, 1.0\}\) weeks), all on the OL+BDC design, using
8-core \texttt{parallel::mclapply}. Wall-clock 483\textasciitilde s;
convergence 100\% across 9\{,\}600 fits.

\end{orig}

At \(c_{bm} = 0.45\):

\begin{table}[htbp]
\centering
\begin{tabular}{lccc}
\toprule
Architecture        & $t_{1/2} = 0$ & $t_{1/2} = 0.5$ & $t_{1/2} = 1.0$ \\
\midrule
MVN (B)             & 0.978         & 0.945           & 0.885 \\
Mean moderation (A) & 0.976         & 0.973           & 0.959 \\
\bottomrule
\end{tabular}
\end{table}

\begin{bullets}

\textbf{The $N = 70$ run confirms the qualitative ordering;
compressed absolute magnitudes reflect ceiling saturation,
not a smaller underlying effect.}

\begin{itemize}\setlength{\itemsep}{2pt}\setlength{\parskip}{0pt}
\item Losses: 9.3 pp (Architecture B, $p < 10^{-12}$) vs.
  1.7 pp (Architecture A, $p \approx 0.05$); ordering matches
  $N = 35$.
\item At $N = 70$ both architectures are within 2--5 pp of the
  ceiling at $t_{1/2} = 0$, mechanically truncating observable
  loss.
\item $N = 35$ production results remain the canonical
  reference for the narrative magnitude claims.
\end{itemize}

\end{bullets}

\begin{rgt}
rgt to complete.

\end{rgt}

\begin{orig}
Within-architecture losses from \(t_{1/2} = 0\) to \(t_{1/2} = 1\) are
9.3 percentage points (Architecture\textasciitilde B; two-proportion
\(z\)-test on the paired cells, \(z = 7.4\), \(p < 10^{-12}\)) and 1.7
percentage points (Architecture\textasciitilde A; \(z = 2.0\),
\(p \approx 0.05\)), where each \(z\)-statistic uses the formula
\((\hat{\pi}_0 - \hat{\pi}_1) /
\sqrt{\hat{\pi}_0(1-\hat{\pi}_0)/n + \hat{\pi}_1(1-\hat{\pi}_1)/n}\)
with \(n = 800\) replicates per cell. The qualitative ordering matches
Section\textasciitilde3.1's \(N = 35\) result. The absolute magnitudes
are compressed because \(N = 70\) at the prazosin-PTSD calibration
places both architectures within 2-5 percentage points of the unit
ceiling at \(t_{1/2} = 0\), mechanically truncating the loss either
architecture's curve can absorb. We note that reading the narrative
loss-magnitude claim against the \(N = 70\) robustness numbers (3\% and
9\% relative loss for Architecture\textasciitilde B versus 0.4\% and
1.8\% for Architecture\textasciitilde A) requires acknowledging this
saturation effect; the \(N = 35\) production results in
Section\textasciitilde3.1 are the canonical reference for the narrative
magnitude.

\end{orig}

\begin{bullets}

\textbf{Type I error is well-controlled across all 18 null
cells, with no architecture-specific inflation.}

\begin{itemize}\setlength{\itemsep}{2pt}\setlength{\parskip}{0pt}
\item Type I error range $[0.010, 0.074]$, mean 0.043, median
  0.042 across 18 null cells at $N = 35$.
\item No architecture-specific inflation detected.
\item The slightly conservative Architecture B OL+BDC null cell
  at $t_{1/2} = 1.0$ ($p = 0.010$) is the only notable
  departure.
\end{itemize}

\end{bullets}

\begin{rgt}
rgt to complete.

\end{rgt}

\begin{orig}
\textbf{Type~I error.} Across all 18 null cells of the \(N = 35\)
production run, Type\textasciitilde I error sat in \([0.010, 0.074]\),
with mean 0.043 and median 0.042. We observe no architecture-specific
inflation. The uniformly sub-nominal rates are consistent with the
\texttt{corCAR1} working-covariance miscalibration characterized in the
companion calibration
study\textsuperscript{\citeproc{ref-pmsimstats-paper10}{1}}; because it
affects both architectures alike, it is common-mode with respect to the
architecture contrast. Architecture\textasciitilde B's \(t_{1/2} = 1.0\)
OL+BDC null (\(p = 0.010\), slightly conservative) is the smallest cell;
the corresponding alternative (\(p = 0.539\)) is the most informative
cell of the Section\textasciitilde3.1 narrative.

\end{orig}

\begin{bullets}

\textbf{Estimator bias is small and architecturally
interpretable; coefficient drift under Architecture A
contrasts with correlation-channel erosion under
Architecture B.}

\begin{itemize}\setlength{\itemsep}{2pt}\setlength{\parskip}{0pt}
\item Mean $\hat\beta_{bm:D_{bc}}$ at $c_{bm} = 0.45$ ranges
  $-0.231$ to $-0.281$ across alternative cells; Architecture A
  OL+BDC at $t_{1/2} = 1.0$ produces the largest absolute
  mean.
\item Architecture B cells maintain mean near $-0.234$ across
  all $t_{1/2}$ values, consistent with erosion through
  correlation decay rather than coefficient drift.
\item Nominal-95\% Wald coverage holds across the grid.
\end{itemize}

\end{bullets}

\begin{rgt}
rgt to complete.

\end{rgt}

\begin{orig}
\textbf{Estimator bias.} Mean estimated \(\hat\beta_{bm{:}D_{bc}}\) at
\(c_{bm} = 0.45\) ranged from \(-0.231\) to \(-0.281\) across the 18
alternative cells, with the \(N = 35\) Architecture-A OL+BDC cell at
\(t_{1/2} = 1.0\) producing the largest absolute mean (\(-0.281\)); the
Architecture\textasciitilde B cells maintained mean near \(-0.234\)
across all \(t_{1/2}\) values, consistent with the Architecture-B
information channel being eroded through correlation decay rather than
through coefficient drift. Empirical SE of \(\hat\beta_{bm{:}D_{bc}}\)
runs 0.05 to 0.10 across cells; the within-replicate model SE matches at
scale, indicating nominal-95\% Wald coverage holds across the grid.

\end{orig}

\bigskip\noindent

\rule{\textwidth}{0.4pt}\bigskip

\section{4. Biological Assumptions}

\begin{bullets}

\textbf{Architecture choice encodes a biological claim about
how the biomarker moderates treatment response.}

\begin{itemize}\setlength{\itemsep}{2pt}\setlength{\parskip}{0pt}
\item The three architectures reflect distinct assumptions about
  the mechanism of biomarker-treatment moderation.
\item Section 4 takes each architecture in turn (4.1, 4.2, 4.3
  for Architectures A, B, and C respectively), then
  discusses the prazosin-PTSD motivating case.
\end{itemize}

\end{bullets}

\begin{rgt}
rgt to complete.

\end{rgt}

\begin{orig}
The choice between architectures is not merely a statistical
convenience; it reflects different assumptions about the biological
mechanism by which the biomarker moderates treatment response. We
consider each architecture in turn, then take up the prazosin-PTSD
motivating application in detail.

\end{orig}

\subsection{4.1 When Architecture A is appropriate}

\begin{bullets}

\textbf{Architecture A is appropriate when the biomarker
directly determines the magnitude of the drug effect; the
scaling is deterministic at the mean-structure level.}

\begin{itemize}\setlength{\itemsep}{2pt}\setlength{\parskip}{0pt}
\item The drug effect scales with the biomarker value:
  $\text{effect}_i = f(B_i) \cdot \text{drug}$.
\item 'Deterministic' means any two participants with the same
  $B_i$ have the same expected response; individual outcomes
  still carry noise.
\item Appropriate when the biomarker is a direct
  pharmacokinetic/pharmacodynamic (PK/PD)
  determinant (eGFR, receptor density, metabolizer status).
\end{itemize}

\end{bullets}

\begin{rgt}
rgt to complete.

\end{rgt}

\begin{orig}
Architecture A (direct mean moderation) is appropriate when the
biomarker governs the \textbf{magnitude of the drug's biological
effect} for each individual participant. The drug effect for participant
\(i\) is scaled by their biomarker value:
\(\text{effect}_i = f(B_i) \cdot \text{drug}\). The scaling is
deterministic in the population-mean sense, in that any two participants
sharing the same \(B_i\) have the same expected response, but individual
outcomes still carry the usual noise, random-effect, and
residual-biological variability. The label `deterministic' refers to the
mean structure of the DGP, not to the realized response of any single
participant.

\end{orig}

\textbf{Clinical examples:}

\begin{itemize}
\item \textbf{Renal clearance of a renally-excreted drug.} Estimated
  glomerular filtration rate (eGFR) directly governs the
  elimination rate of drugs such as digoxin or aminoglycoside
  antibiotics. Participants with higher eGFR clear the drug more
  rapidly, achieving lower steady-state concentrations and a
  proportionally smaller drug effect. The biomarker-drug
  relationship is mechanistic and deterministic: given the
  biomarker value, the effective dose at steady state is fully
  determined (up to random noise).

\item \textbf{Receptor density moderation.} A biomarker that measures
  target receptor density (e.g., PET imaging of receptor
  availability for an antagonist) directly determines the
  pharmacodynamic response to that antagonist. More receptors
  means more drug effect, proportionally.

\item \textbf{Genetic metabolizer status.} CYP2D6 metabolizer phenotype
  determines the rate of drug activation or clearance. Poor
  metabolizers of a prodrug receive less active compound,
  reducing their treatment effect proportionally.

\item \textbf{Dose-response with biomarker-dependent effective dose.}
  When the biomarker determines the effective dose (through
  body composition, protein binding, or volume of distribution),
  the treatment effect scales with the biomarker through the
  dose-response curve.
\end{itemize}

\begin{bullets}

\textbf{In all Architecture A cases, the proportional
biomarker-drug relationship is preserved during the off-drug
period.}

\begin{itemize}\setlength{\itemsep}{2pt}\setlength{\parskip}{0pt}
\item Residual drug effect during the off-drug period maintains the same
  proportional relationship with the biomarker.
\item If participant A has 2$\times$ the drug effect of B
  on-drug, they also have $\approx$2$\times$ the residual
  effect during the off-drug period.
\item This preserved proportionality is what Architecture A
  encodes and why it loses little power under carryover.
\end{itemize}

\end{bullets}

\begin{rgt}
rgt to complete.

\end{rgt}

\begin{orig}
In all these cases the biomarker's moderating effect is preserved under
carryover because the residual drug effect during off-drug periods
maintains the same proportional relationship with the biomarker. If
participant A has twice the drug effect of participant B when on drug,
they will also have approximately twice the residual effect during the
off-drug period. That proportionality is precisely what Architecture A
encodes.

\end{orig}

\subsection{4.2 When Architecture B is appropriate}

\begin{bullets}

\textbf{Architecture B is appropriate when the biomarker is a
probabilistic predictor of response, mediated by latent
factors the biomarker imperfectly indexes.}

\begin{itemize}\setlength{\itemsep}{2pt}\setlength{\parskip}{0pt}
\item Appropriate when the biomarker predicts which participants
  are likely to respond, not the magnitude of any individual's
  response.
\item The biomarker-response association is population-level and
  probabilistic, mediated by latent factors (disease subtype,
  neurobiological heterogeneity, comorbidities).
\item Two participants with identical biomarker values may
  respond very differently.
\end{itemize}

\end{bullets}

\begin{rgt}
rgt to complete.

\end{rgt}

\begin{orig}
Architecture B (differential correlation) is appropriate when the
biomarker predicts \textbf{which participants are likely to
respond}, but the magnitude of any individual's response is not
deterministically governed by the biomarker. The biomarker and the drug
response are associated at the population level, but the association is
mediated by latent factors (disease subtype, neurobiological
heterogeneity, comorbidities) that the biomarker imperfectly indexes.

\end{orig}

\textbf{Clinical examples:}

\begin{itemize}
\item
  \textbf{Blood pressure as a PTSD subtype marker.} Elevated resting
  blood pressure may index a noradrenergic-predominant PTSD
  phenotype\textsuperscript{\citeproc{ref-hendrickson2020}{13}} that is
  more likely to respond to prazosin (an alpha-1 adrenergic antagonist).
  The blood pressure does not \textit{cause} the drug to work better;
  rather, it correlates with an underlying neurobiological state that
  determines drug responsiveness. Two participants with identical blood
  pressure may have very different responses because blood pressure is
  an imperfect proxy for the noradrenergic state.
\item
  \textbf{Baseline disease severity as a treatment predictor.} Trials in
  Alzheimer's disease, depression, and chronic pain often find that
  patients with more severe baseline symptoms show larger treatment
  effects (the `floor effect' or `regression to the mean' phenomenon).
  The baseline severity score correlates with treatment response, but
  the relationship is statistical (population-level tendency), not
  deterministic.
\item
  \textbf{Inflammatory biomarkers in psychiatry.} C-reactive protein
  (CRP) or interleukin-6 levels may predict response to
  anti-inflammatory augmentation of
  antidepressants\textsuperscript{\citeproc{ref-raison2013}{18}}.
  High-CRP patients are more likely to respond, but the relationship is
  probabilistic: many high-CRP patients do not respond, and some low-CRP
  patients do.
\item
  \textbf{Genetic risk scores.} Polygenic risk scores predict disease
  trajectory and treatment response, but explain only a fraction of the
  variance. The remainder reflects unmeasured genetic, epigenetic, and
  environmental factors.
\end{itemize}

\begin{bullets}

\textbf{Under Architecture B, carryover weakens the
biomarker-response correlation because the biomarker's
predictive power depends on the presence of active drug.}

\begin{itemize}\setlength{\itemsep}{2pt}\setlength{\parskip}{0pt}
\item As the drug washes out, the drug-mediated variance
  component diminishes, weakening the BM-BR correlation.
\item The biomarker's predictive power is tied to active drug
  effect; without drug, it predicts nothing about
  non-drug symptom change.
\item This is the 'correlation erosion' channel unique to
  Architecture B.
\end{itemize}

\end{bullets}

\begin{rgt}
rgt to complete.

\end{rgt}

\begin{orig}
In these cases carryover has a different implication: as the drug washes
out, the correlation between biomarker and response weakens because the
drug-mediated component of response variance diminishes. The biomarker's
predictive power is inherently tied to the presence of active drug
effect. Without drug, the biomarker is just a biomarker; it does not
predict the non-drug component of symptom change.

\end{orig}

\begin{bullets}

\textbf{Architecture B is a tractable second-moment
approximation to the biologically faithful finite-mixture DGP;
adequate for power calculations but not a mechanistic
endorsement.}

\begin{itemize}\setlength{\itemsep}{2pt}\setlength{\parskip}{0pt}
\item If the biomarker indexes a discrete responder class, the
  faithful DGP is a finite mixture; Architecture B approximates
  its population-level covariance structure.
\item Architecture B captures the covariance implications of
  latent-class membership without committing to the full
  generative form.
\item Adequate for power calculations; should not be read as
  endorsing MVN differential correlation as the correct
  mechanism.
\end{itemize}

\end{bullets}

\begin{rgt}
rgt to complete.

\end{rgt}

\begin{orig}
The latent-subtype framing raises a legitimate modeling question: if the
biomarker is really a noisy indicator of an underlying responder class,
the biologically faithful DGP would be an explicit finite mixture in
which each participant is drawn from one of two (or more) response
distributions with class-membership probability a function of \(B_i\).
Under such a mixture DGP the individual-level moderation is discrete,
not probabilistic, and the MVN differential-correlation parameterization
is best understood as a tractable second-moment approximation to a
richer mixture model rather than as the generative mechanism itself. The
practical consequence is that Architecture\textasciitilde B as
implemented here captures the population-level covariance implications
of a latent-class mechanism without committing to its full generative
form; this is adequate for power calculations that test for the presence
of an interaction but should not be read as an endorsement of MVN
differential correlation as the correct mechanistic model. Whether the
observed power loss under carryover is better mitigated by analysis
strategies (Section 6) or by explicit mixture modeling is taken up in a
companion paper\textsuperscript{\citeproc{ref-pmsimstats-paper03}{14}},
which derives the conditions under which a single multivariate normal
reproduces the first two moments of a two-component mixture, catalogues
the regimes in which that approximation fails, and finds in a
preliminary evaluation that at trial-relevant sample sizes class-aware
analyses are not competitive with the continuous linear-mixed baseline.

\end{orig}

\begin{bullets}

\textbf{Senn (2016) cautions that apparent personalized
response is often a variance artifact; Architecture B
presupposes a real, stable latent subtype, which may not hold.}

\begin{itemize}\setlength{\itemsep}{2pt}\setlength{\parskip}{0pt}
\item Senn (2016) argues that apparent biomarker-treatment
  interactions are often variance-component artifacts, not
  genuine effects.
\item Architecture B presupposes the latent subtype is real,
  identifiable, and stable across the trial period.
\item If that presupposition is shaky, the documented power
  loss may overstate recoverable signal.
\end{itemize}

\end{bullets}

\begin{rgt}
rgt to complete.

\end{rgt}

\begin{orig}
The latent-subtype framing of Architecture B should also be adopted with
appropriate epistemic caution.
Senn\textsuperscript{\citeproc{ref-senn2016mastering}{19}} has argued
that the apparent `personal element in response to treatment' is often a
variance-component artifact rather than a genuine biomarker-treatment
interaction, and that even modest within-subject variability can produce
convincing-looking but spurious differential response. Architecture B
presupposes that the latent subtype indexed by the biomarker is real,
identifiable, and stable across the trial period. When that
presupposition is shaky, the apparent power loss under carryover
documented in Section 3 may overstate what could be recovered under any
analysis model, because the underlying signal it is attempting to
estimate may itself be smaller than assumed.

\end{orig}

\subsection{4.3 The prazosin-PTSD case}

We turn briefly to the motivating clinical application for pmsimstats,
which is the active-duty prazosin trial by Raskind et
al.\textsuperscript{\citeproc{ref-raskind2013}{15}}, whose
SBP-stratified secondary analysis is the subject of ongoing reanalysis
by the pmsimstats team. The biomarker is resting systolic blood pressure
(SBP) and the outcome is CAPS (Clinician-Administered PTSD Scale) score
change.

\begin{bullets}

\textbf{The prazosin-PTSD evidence base is mixed; the PACT
null result motivates the search for a predictive biomarker,
with SBP as the leading candidate.}

\begin{itemize}\setlength{\itemsep}{2pt}\setlength{\parskip}{0pt}
\item Earlier single-site Raskind trials (2003, 2007, 2013) reported benefit on nightmares and
  PTSD symptoms.
\item PACT (2018) failed to replicate; this has been
  interpreted as consistent with effect heterogeneity.
\item SBP, as a noninvasive proxy for central noradrenergic
  tone, is the leading predictive biomarker candidate.
\end{itemize}

\end{bullets}

\begin{rgt}
rgt to complete.

\end{rgt}

\begin{orig}
The clinical evidence base for prazosin in PTSD is mixed. The earlier
single-site trials by Raskind and
colleagues\textsuperscript{\citeproc{ref-raskind2013}{15},\citeproc{ref-raskind2003}{20},\citeproc{ref-raskind2007}{21}}
reported benefit on nightmares and global PTSD symptoms in veteran and
active-duty samples. The subsequent multisite Veterans Affairs
Cooperative Study Program
trial\textsuperscript{\citeproc{ref-raskind2018}{16}}, known as PACT,
did not replicate the effect: prazosin failed to outperform placebo on
the primary endpoints. The negative PACT result has been interpreted by
several investigators as consistent with effect heterogeneity, namely
that prazosin works in some patients but not others, which motivates the
search for a predictive biomarker that would identify the responder
subgroup. Resting SBP, as a noninvasive proxy for central noradrenergic
tone, is the leading candidate.

\end{orig}

\begin{bullets}

\textbf{The SBP-prazosin mechanism is an Architecture B
scenario: SBP is a proxy for noradrenergic state, not a direct
PK determinant.}

\begin{itemize}\setlength{\itemsep}{2pt}\setlength{\parskip}{0pt}
\item Elevated SBP indexes heightened noradrenergic tone, a
  PTSD subtype preferentially responsive to prazosin.
\item SBP is a proxy, not a causal determinant; two patients
  with identical SBP may have very different noradrenergic
  states.
\item This is precisely the Architecture B scenario:
  probabilistic association, not deterministic scaling.
\end{itemize}

\end{bullets}

\begin{rgt}
rgt to complete.

\end{rgt}

\begin{orig}
The biological hypothesis is that elevated SBP indexes heightened
central noradrenergic tone, which characterizes a subtype of PTSD that
is preferentially responsive to alpha-1 adrenergic blockade by prazosin.
This is an Architecture B scenario: SBP is a proxy for an underlying
neurobiological state, not a direct determinant of drug
pharmacokinetics. Two patients with the same SBP may differ in their
actual noradrenergic tone, alpha-1 receptor distribution, and downstream
signaling, leading to different treatment responses despite identical
biomarker values.

\end{orig}

\begin{bullets}

\textbf{Hendrickson et al. (2020) correctly uses Architecture B
for this context; the carryover power loss is clinically
relevant for trial design decisions.}

\begin{itemize}\setlength{\itemsep}{2pt}\setlength{\parskip}{0pt}
\item Architecture B (MVN differential correlation) is
  appropriate for the SBP-prazosin mechanism.
\item The documented carryover sensitivity means designs with
  longer off-drug periods will have meaningfully less power.
\item This effect is larger than a direct mean moderation
  simulation would predict.
\end{itemize}

\end{bullets}

\begin{rgt}
rgt to complete.

\end{rgt}

\begin{orig}
The Hendrickson et
al.\textsuperscript{\citeproc{ref-hendrickson2020}{13}} DGP uses
Architecture B (MVN differential correlation) for this clinical context.
The finding that carryover substantially reduces power under this
architecture is therefore clinically relevant: it suggests that trial
designs with longer off-drug periods will have meaningfully less power
to detect the SBP-prazosin interaction than designs with shorter
off-drug periods, and this effect is larger than what a direct mean
moderation simulation would predict.

\end{orig}

\bigskip\noindent

\rule{\textwidth}{0.4pt}\bigskip

\section{5. Implications for Simulation Study Design}

\subsection{5.1 Choosing the appropriate architecture}

The choice of DGP architecture should be guided by the biological
mechanism hypothesized for the biomarker-treatment interaction:

\begin{table}[htbp]
\centering
\begin{tabular}{lcc}
\toprule
Criterion & Architecture A & Architecture B \\
\midrule
Biomarker role   & Causal mediator       & Statistical predictor \\
Mechanism        & Deterministic scaling & Probabilistic association \\
Signal location  & Mean structure        & Covariance structure \\
Free parameters  & 1 ($c_{bm}$)          & 1 ($c_{bm}$) \\
Carryover impact & Modest (1--3\%)       & Design-dependent, up to 31\% \\
Appropriate when &
  \parbox[t]{3cm}{Biomarker\\determines dose\\or PK} &
  \parbox[t]{3cm}{Biomarker\\indexes a latent\\subtype} \\
\bottomrule
\end{tabular}
\end{table}

When the operative mechanism is unknown, Architecture\textasciitilde B
is the conservative default for trial-design decisions: its larger
carryover sensitivity keeps power estimates from being optimistic
relative to Architecture\textasciitilde A. A combined architecture that
mixes the two channels, together with a corresponding sensitivity sweep,
is developed in the companion
paper\textsuperscript{\citeproc{ref-pmsimstats-paper11}{17}}.

\subsection{5.2 Reporting requirements}

Simulation studies evaluating biomarker-moderated treatment effects
should explicitly state:

\begin{enumerate}
\item Whether the biomarker-treatment interaction is generated through
  mean moderation, differential correlation, or both channels
  simultaneously.
\item The biological rationale for the chosen mechanism.
\item Whether the carryover model (if present) operates on the
  interaction signal consistently with the chosen architecture.
\item Sensitivity analyses under alternative architectures, if the
  biological mechanism is uncertain.
\end{enumerate}

\begin{bullets}

\textbf{The four reporting requirements are the ADEMP 'D'
element applied specifically to biomarker-validation
simulations.}

\begin{itemize}\setlength{\itemsep}{2pt}\setlength{\parskip}{0pt}
\item ADEMP (2019) prescribes five elements:
  Aims, Data-generating mechanisms, Estimands, Methods,
  Performance measures.
\item The DGP-architecture choice concerns the D element;
  the four requirements above specialize that element for
  architecture-sensitive simulations.
\end{itemize}

\end{bullets}

\begin{rgt}
rgt to complete.

\end{rgt}

\begin{orig}
The reporting framework above is consistent with the broader ADEMP
template of Morris, White, and
Crowther\textsuperscript{\citeproc{ref-morris2019simulation}{22}}, which
prescribes that simulation studies fix five elements before production
runs: the \textbf{Aims} of the study, the
\textbf{Data-generating mechanisms}, the \textbf{Estimands} (the
inferential targets), the \textbf{Methods} (analysis specifications
compared), and the \textbf{Performance measures} (power, bias, coverage,
Type I error, and MCSEs of each). The DGP-architecture choice this paper
concerns is the \textbf{D} element of an ADEMP specification; the
reporting requirements above are the specialization of the broader ADEMP
framework to the architecture-sensitive class of biomarker-validation
simulations.

\end{orig}

\begin{bullets}

\textbf{Primary estimand: $\beta_{bm:D}$ in the linear
mixed-effects (LME) model;
performance measures cover power, Type I error, bias, SE, and
convergence; 1000 replicates per cell.}

\begin{itemize}\setlength{\itemsep}{2pt}\setlength{\parskip}{0pt}
\item Estimand: $\beta_{bm:D}$ in
  \texttt{Sx \textasciitilde{} bm + t + Dbc + bm:Dbc} with random
  intercept and AR(1).
\item Five performance measures: power, Type I error, bias,
  empirical SE, convergence rate.
\item $n_{\text{reps}} = 1000$ throughout; calibrated to
  achieve MCSE $< 0.02$ for power near 0.5.
\end{itemize}

\end{bullets}

\begin{rgt}
rgt to complete.

\end{rgt}

\begin{orig}
The \textbf{primary estimand} for this paper's simulations is the
biomarker-treatment interaction coefficient \(\beta_{bm:D}\) in the
linear-mixed analysis model
\texttt{Sx \textasciitilde{} bm + t + Dbc + bm:Dbc} with random
intercept and AR(1) residual correlation. The reported
\textbf{performance measures} are: (i) power for \(\beta_{bm:D}\) at
\(\alpha = 0.05\), with MCSE
\(\sqrt{\hat{\pi}(1-\hat{\pi})/n_{\text{reps}}}\); (ii) Type I error
under \(c_{bm} = 0\), with the same MCSE; (iii) bias of
\(\hat{\beta}_{bm:D}\) relative to the calibrated true value at each
parameter cell; (iv) the empirical MCSE of \(\hat{\beta}_{bm:D}\) across
replicates; and (v) the convergence rate of the linear-mixed fit. Sample
size, biomarker effect size, and carryover half-life vary across the
parameter grid; the number of replicates per cell is calibrated to
achieve MCSE below 0.02 for power at \(\hat{\pi} = 0.5\), which
corresponds to \(n_{\text{reps}} \geq 625\). The simulations reported
here use \(n_{\text{reps}} = 1000\) throughout.

\end{orig}

\subsection{5.3 When the choice is uncertain}

If the biological mechanism is ambiguous, as it often is in early
biomarker development, we suggest investigators take the following
steps:

\begin{enumerate}
\item Run the simulation under both architectures and report the
  range of power estimates.
\item Architecture B produces more conservative (lower)
  power estimates under carryover, making it the safer default
  for trial design decisions.
\item Design the trial to minimize carryover (adequate washout
  periods) to reduce sensitivity to the DGP architecture choice.
\end{enumerate}

\bigskip\noindent

\rule{\textwidth}{0.4pt}\bigskip

\section{6. Alternative Analysis Strategies Under Architecture B}

\emph{This section provides a conceptual overview of six
candidate strategies for mitigating the power loss documented in
Section 3. None of the strategies below has been evaluated
by simulation within this study; the assessments of expected
benefit are qualitative and based on the mechanistic account
in Section 3.2. A systematic simulation comparison is warranted
before any strategy is recommended for applied use.}

The power loss documented in Section 3 arises from two distinct sources:
(1) compressed variance of the drug exposure variable \texttt{Dbc} when
carryover is present, and (2) erosion of the BM-BR correlation during
off-drug periods. Source 1 is a property of the analysis model
parameterization and can potentially be addressed by choosing a
different predictor. Source 2 is a property of the data itself (the
signal is genuinely weaker in off-drug observations), and no analysis
model can recover signal that is not present. However, analysis
strategies that are selective about \textit{which observations
contribute to the test} or \textit{how they are weighted} can mitigate
the damage.

\subsection{6.1 Restrict analysis to on-drug observations}

The most direct approach is to discard off-drug timepoints entirely and
test the biomarker-treatment interaction using only on-drug
observations. Every retained observation has the full BM-BR correlation
(\(c_{bm}\)), so correlation erosion is eliminated. The costs are
twofold: a reduction in effective sample size (fewer observations per
participant) and, more importantly, the loss of the off-drug reference
that separates the biological response from placebo and time-varying
background response. Without off-drug observations the biological,
placebo, and time-varying components are confounded, which is why the
Open-Label design, despite having the most on-drug timepoints, does not
dominate the other designs in the Section 3 results. For designs with
many on-drug timepoints (OL has 8, the Hybrid design has 5-6 per path),
the net effect may be positive: the per-observation signal-to-noise
ratio is maximal, and the analysis model simplifies to
\texttt{Sx \textasciitilde{} bm + t + bm:t} since there is no treatment
contrast to model. The interaction is detected through the correlation
between biomarker and the time-on-drug trajectory.

\begin{bullets}

\textbf{Restricting to on-drug observations eliminates
correlation erosion but loses the off-drug reference; best
for OL and Hybrid designs, worst for CO.}

\begin{itemize}\setlength{\itemsep}{2pt}\setlength{\parskip}{0pt}
\item Benefit: every retained observation has full $c_{bm}$;
  correlation erosion is eliminated.
\item Cost: loss of off-drug reference confounds biological,
  placebo, and time-varying components.
\item Practical: best for OL/Hybrid (many on-drug timepoints);
  least useful for CO (discards half the observations).
\end{itemize}

\end{bullets}

\begin{rgt}
rgt to complete.

\end{rgt}

\begin{orig}
This approach is most promising for the Open-Label design (which has no
off-drug observations regardless) and the Hybrid design (which has a
long initial on-drug phase). It is least useful for the Crossover
design, where discarding half the observations substantially reduces
statistical power.

\end{orig}

\subsection{6.2 Weighted analysis}

\begin{bullets}

\textbf{Weighted analysis downweights off-drug observations by
drug exposure without discarding them; information-theoretic
rationale; straightforward to implement.}

\begin{itemize}\setlength{\itemsep}{2pt}\setlength{\parskip}{0pt}
\item Weight: $w_t = e^{-\lambda t_{sd}}$ for off-drug
  observations; 1 for on-drug.
\item Rationale: higher exposure = stronger correlation signal
  = more information about the interaction.
\item Implemented via the \texttt{weights} argument in
  \texttt{nlme::lme}.
\end{itemize}

\end{bullets}

\begin{rgt}
rgt to complete.

\end{rgt}

\begin{orig}
Rather than discarding off-drug observations entirely, a weighted
analysis retains all observations but weights them by their estimated
drug exposure level. On-drug observations receive weight 1. Off-drug
observations receive weight proportional to the expected residual drug
effect:

\[w_t = e^{-\lambda \cdot t_{sd}}\]

This downweights observations where the BM-BR correlation has decayed
(and where they contribute more noise than signal to the interaction
estimate) without discarding them entirely. The rationale is
information-theoretic: observations with higher drug exposure carry more
information about the biomarker-treatment interaction because the
correlation signal is stronger. This is straightforward to implement in
\texttt{nlme::lme} using the \texttt{weights} argument.

\end{orig}

\subsection{6.3 Within-subject contrast}

For designs with both on-drug and off-drug periods (CO, Hybrid, OL+BDC),
a summary-measure approach computes a per-participant contrast: the
difference between mean on-drug response and mean off-drug response. The
biomarker-treatment interaction is then tested by a simple regression of
these contrasts on the biomarker:

\[\bar{Y}_{i,\text{on}} - \bar{Y}_{i,\text{off}} = \alpha +
\gamma \cdot B_i + \epsilon_i\]

\begin{bullets}

\textbf{Within-subject contrast sidesteps the repeated-measures
model; trades efficiency for robustness to carryover
misspecification.}

\begin{itemize}\setlength{\itemsep}{2pt}\setlength{\parskip}{0pt}
\item Per-participant contrast $\bar{Y}_{\text{on}} -
  \bar{Y}_{\text{off}}$ is regressed on the biomarker.
\item Carryover compresses the contrast magnitude but the
  contrast-biomarker correlation may be more robust.
\item Trade-off: less efficient than a full longitudinal model,
  but more robust to carryover and correlation misspecification.
\end{itemize}

\end{bullets}

\begin{rgt}
rgt to complete.

\end{rgt}

\begin{orig}
This sidesteps the \texttt{Dbc} parameterization and the
repeated-measures model entirely. Carryover affects the magnitude of the
contrast (making off-drug means closer to on-drug means), but the
\textit{correlation between the contrast
and the biomarker} may be more robust because within-subject averaging
reduces noise before the biomarker association is tested. The approach
trades the efficiency of a full longitudinal model for robustness to
misspecification of the carryover and correlation structure.

\end{orig}

\subsection{6.4 Two-stage random slopes}

A two-stage approach first estimates participant-specific treatment
effects, then tests their association with the biomarker. Stage 1 fits a
random-slopes model:

\[Y_{it} = \beta_0 + \beta_1 t + \beta_2 D_{it} + u_{0i} +
u_{1i} D_{it} + \epsilon_{it}\]

where \(\beta_2\) is the population mean drug effect and \(u_{1i}\) is
participant \(i\)'s random deviation from that mean. The
participant-specific drug effect is then \(\beta_2 + u_{1i}\). Stage 2
regresses the estimated random deviations (or equivalently, the best
linear unbiased predictors (BLUPs) of the participant-specific effects)
on the biomarker:

\[\hat{u}_{1i} = \delta_0 + \delta_1 B_i + \eta_i\]

\begin{bullets}

\textbf{Two-stage random slopes separates within-subject effect
estimation from between-subject biomarker association; may be
more robust to carryover.}

\begin{itemize}\setlength{\itemsep}{2pt}\setlength{\parskip}{0pt}
\item Stage 1: random-slopes LME estimates participant-specific
  drug effects ($\beta_2 + u_{1i}$).
\item Stage 2: regress BLUPs on biomarker to test association.
\item The random slope $u_{1i}$ integrates over the full
  trajectory; robustness to carryover is plausible but
  unverified by simulation.
\end{itemize}

\end{bullets}

\begin{rgt}
rgt to complete.

\end{rgt}

\begin{orig}
This separates the within-subject treatment effect estimation (which
uses the longitudinal data and is affected by carryover) from the
between-subject biomarker association (which uses the cross-sectional
relationship between estimated effects and biomarker values). The
approach may be more robust to carryover because the random slope
\(u_{1i}\) integrates over the participant's entire trajectory,
including both on-drug and off-drug phases.

\end{orig}

\subsection{6.5 Exclusion of contaminated observations}

\begin{bullets}

\textbf{Excluding the first 1--2 off-drug timepoints removes
the most contaminated observations; later off-drug timepoints
provide a cleaner contrast.}

\begin{itemize}\setlength{\itemsep}{2pt}\setlength{\parskip}{0pt}
\item Early off-drug timepoints: highest residual drug effect,
  BM-BR correlation actively decaying.
\item Later off-drug timepoints: fully washed out, cleaner
  contrast with on-drug phase.
\item Result: comparison between fully-on and fully-off states,
  avoiding the strained intermediate zone.
\end{itemize}

\end{bullets}

\begin{rgt}
rgt to complete.

\end{rgt}

\begin{orig}
The first 1-2 off-drug timepoints carry the most carryover
contamination: the residual drug effect is highest, and the BM-BR
correlation is in the process of decaying. Later off-drug timepoints,
where the drug has fully washed out, provide a cleaner contrast with the
on-drug phase. Excluding the early off-drug observations and retaining
only the later ones yields a comparison between fully-on and fully-off
states, avoiding the intermediate zone where the analysis model is most
strained.

\end{orig}

\subsection{6.6 Design-level solutions}

Rather than adapting the analysis model to accommodate carryover, the
trial design itself can be modified to reduce carryover exposure:

\begin{itemize}
\item \textbf{Longer washout periods} between drug phases eliminate
  residual drug effect before off-drug measurements begin.
\item \textbf{More on-drug timepoints} relative to off-drug increase the
  proportion of high-signal observations.
\item \textbf{Front-loaded drug exposure} (as in the OL+BDC design)
  places a long initial on-drug block before the off-drug
  baseline-during-continued-treatment phase, so that the BM-BR
  correlation has been fully established before any off-drug
  observations are collected. The alternative ordering
  (off-drug first, then on-drug) has the appealing property of
  placing the placebo and time-varying estimation in an
  uncontaminated window but sacrifices on-drug exposure duration
  and forces the BM-BR signal to be re-established after a
  wash-in; the net power trade-off depends on the ratio of
  carryover half-life to total trial length and is explored
  empirically in the Section 3 results.
\end{itemize}

\begin{bullets}

\textbf{Design-level changes address the root cause rather than
the symptom; cost is longer trial or reduced ability to estimate
carryover dynamics.}

\begin{itemize}\setlength{\itemsep}{2pt}\setlength{\parskip}{0pt}
\item Three levers: longer washout, more on-drug timepoints,
  front-loaded drug exposure (as in OL+BDC).
\item Front-loading establishes BM-BR correlation before any
  off-drug observations; alternative ordering (off-drug first)
  trades clean placebo estimation for reduced on-drug exposure.
\item Addresses root cause; cost is typically longer trial or
  reduced carryover estimation ability.
\end{itemize}

\end{bullets}

\begin{rgt}
rgt to complete.

\end{rgt}

\begin{orig}
Design-level solutions address the root cause (carryover in the data)
rather than the symptom (power loss in the analysis). Their cost is
typically a longer trial duration or reduced ability to estimate the
carryover dynamics themselves.

\end{orig}

\subsection{6.7 Qualitative comparison}

Table\textasciitilde{}\ref{tab:sec6-strategies} summarizes the
strategies along three dimensions: whether they discard observations,
implementation complexity, and the expected direction of benefit based
on the mechanistic account in Section 3.2. All entries in the
`Hypothesized benefit' column are qualitative; no simulation evidence is
presented here, and each strategy should be evaluated empirically before
use in trial planning.

\begin{table}[htbp]
\centering
\caption{Qualitative summary of analysis strategies for
recovering power under Architecture B carryover. All benefit
assessments are mechanistic hypotheses, not simulation-supported
estimates.}
\label{tab:sec6-strategies}
\begin{tabular}{lccp{5cm}}
\toprule
Strategy & Discards data & Complexity & Hypothesized benefit \\
\midrule
On-drug only & Yes & Low & Eliminates correlation erosion;
  forfeits off-drug reference (best for OL/Hybrid) \\
Weighted & No & Low & Reduces erosion damage proportionally;
  depends on weight calibration \\
Within-subject contrast & Aggregates & Low & Robust to
  carryover misspecification; less efficient than LME \\
Two-stage random slopes & No & Moderate & Unknown; separates
  estimation stages; robustness unverified \\
Exclude early off-drug & Some & Low & Removes most contaminated
  observations; improves contrast clarity \\
Design modification & N/A & N/A & Addresses root cause;
  cost is longer trial or reduced carryover estimation \\
\bottomrule
\end{tabular}
\end{table}

A systematic simulation comparison of these strategies under
Architecture B, analogous to the factorial comparison of analysis models
in Section 3, would identify which approach best recovers the power lost
to carryover for each trial design. We regard this comparison as
important future work and note that the on-drug-only restriction and the
weighted analysis are the most tractable candidates for a first
empirical evaluation.

\bigskip\noindent

\rule{\textwidth}{0.4pt}\bigskip

\section{7. Relationship to the Broader Literature}

\subsection{7.1 Treatment effect heterogeneity}

\begin{bullets}

\textbf{The PATH Statement distinguishes risk modeling from effect
modeling for predictive HTE analysis; this is an analysis-model
distinction, not a DGP distinction.}

\begin{itemize}\setlength{\itemsep}{2pt}\setlength{\parskip}{0pt}
\item Rothwell (2005) establishes the predictive/prognostic
  biomarker distinction; Kent et al. (2020) synthesize it into
  the PATH Statement.
\item PATH's risk modeling approach stratifies on a multivariable
  risk score; effect modeling adds explicit interaction terms.
\item PATH favors risk modeling and cautions against data-driven
  effect modeling because of low power, multiplicity, and
  overfitting.
\end{itemize}

\end{bullets}

\begin{rgt}
rgt to complete.

\end{rgt}

\begin{orig}
The statistical literature on treatment effect heterogeneity is
extensive. Rothwell\textsuperscript{\citeproc{ref-rothwell2005}{3}} sets
out the foundational distinction between \textit{predictive} biomarkers
(those that modify the treatment effect) and \textit{prognostic}
biomarkers (those that predict outcome regardless of treatment). The
current consensus statement on predictive HTE analysis is the Predictive
Approaches to Treatment effect Heterogeneity (PATH) Statement of Kent et
al.\textsuperscript{\citeproc{ref-kent2020path}{4}}, which distinguishes
two analytic strategies for identifying predictive biomarker effects in
randomized trials: a \textit{risk modeling} approach, in which a
multivariable risk model (developed without a treatment term) is used to
stratify patients and examine variation in absolute treatment benefit
across risk strata, and an \textit{effect modeling} approach, in which
the regression includes explicit treatment-by-covariate interaction
terms. PATH expresses caution about data-driven effect modeling on the
grounds of low statistical power, multiplicity, and overfitting, and
favors risk modeling when the application criteria are met.

\end{orig}

\begin{bullets}

\textbf{PATH's analysis-model axis (risk vs.\ effect modeling)
is orthogonal to our DGP axis; both PATH strategies implicitly
assume Architecture A.}

\begin{itemize}\setlength{\itemsep}{2pt}\setlength{\parskip}{0pt}
\item PATH's distinction is analysis-model; ours is DGP; they
  are orthogonal.
\item Both PATH strategies assume heterogeneity is in the mean
  structure (Architecture A); Architecture B does not fit
  either PATH category.
\item Architecture A = parametric interaction model;
  Architecture B = latent-class/mixture-model perspective.
\end{itemize}

\end{bullets}

\begin{rgt}
rgt to complete.

\end{rgt}

\begin{orig}
We note that the PATH risk-versus-effect-modeling distinction is an
analysis-model axis. The Architecture A versus Architecture B
distinction introduced in this paper is a data-generating-process axis,
and the two are orthogonal. Both PATH analytic strategies implicitly
assume that the underlying treatment effect heterogeneity is encoded in
the mean structure of the outcome (the heterogeneity is identifiable
from differences in conditional means), which corresponds to
Architecture A as a DGP. Architecture B as a DGP is a separate
phenomenon that does not map cleanly onto either of PATH's
analysis-model categories: the interaction signal resides in the
covariance structure rather than the mean, and the standard
regression-based predictive HTE analysis methods endorsed by PATH may
have reduced power under Architecture B precisely because they were
developed under an implicit Architecture A assumption. Architecture A
can accordingly be characterized as a
\textit{parametric interaction model} (in PATH's effect modeling sense).
Architecture B can be characterized as a \textit{latent
class} or \textit{mixture model} perspective in which the biomarker is a
noisy indicator of class membership.

\end{orig}

\subsection{7.2 Prevalence of each architecture}

\begin{bullets}

\textbf{Architecture A is the near-universal standard across the
biomarker-moderated trial simulation literature; Architecture B
is absent from all major threads except Hendrickson et al.}

\begin{itemize}\setlength{\itemsep}{2pt}\setlength{\parskip}{0pt}
\item Simon (2010), Freidlin and Korn (2014), Renfro and Sargent
  (2017), Haller and Ulm (2018), and the platform trial
  literature all use Architecture A.
\item N-of-1 simulation studies (1997, 2016, 2013) use hierarchical random-effects
  models, a variant of Architecture A.
\item Wang and Schork (2019) and Schork (2022) share authorship
  with Hendrickson et al. yet adopt Architecture A, suggesting
  the differential-correlation choice is paper-specific.
\end{itemize}

\end{bullets}

\begin{rgt}
rgt to complete.

\end{rgt}

\begin{orig}
A survey of the simulation methodology literature reveals that
Architecture A (direct mean moderation) is the near-universal standard.
Simon\textsuperscript{\citeproc{ref-simon2010}{5}} employs it throughout
the simulation frameworks he develops for biomarker-stratified designs;
Freidlin and Korn\textsuperscript{\citeproc{ref-freidlin2014}{6}} adopt
it for adaptive enrichment trials; Renfro and
Sargent\textsuperscript{\citeproc{ref-renfro2017}{7}} use it for basket
and umbrella trials; and it underlies the platform trial literature
throughout. Haller and
Ulm\textsuperscript{\citeproc{ref-haller2018}{8}}, in simulation work on
biomarker-treatment interaction estimation in randomized trials with
prognostic covariates, likewise adopt the standard regression
specification, focusing on power and bias of the interaction estimator
under different covariate adjustment strategies. N-of-1 simulation
studies, including Zucker et al. (1997), Araujo et
al.\textsuperscript{\citeproc{ref-araujo2016}{10}}, and Duan et
al.\textsuperscript{\citeproc{ref-duan2013}{23}}, use hierarchical
random-effects models in which individual treatment effects are drawn
from a population distribution, a variant of Architecture A in which the
biomarker predicts the random treatment effect through the mean
structure. Methodological work on power and design for N-of-1 trials
with serial correlation and carryover, from both Wang and
Schork\textsuperscript{\citeproc{ref-wang2019}{11}} and
Schork\textsuperscript{\citeproc{ref-schork2022}{12}}, adopts the same
Architecture A specification, even though this work shares senior
authorship with Hendrickson et
al.\textsuperscript{\citeproc{ref-hendrickson2020}{13}}. It appears,
then, that the differential- correlation parameterization in the latter
is a paper-specific design choice rather than a methodological
convention of any laboratory or research program.

\end{orig}

\begin{bullets}

\textbf{Architecture B appears in latent variable / structural
equation modeling (SEM) /
copula contexts but is unique in N-of-1 clinical trial
simulation.}

\begin{itemize}\setlength{\itemsep}{2pt}\setlength{\parskip}{0pt}
\item Architecture B appears in latent variable / SEM /
  Bayesian joint modeling / copula contexts (e.g., Rizopoulos
  2012).
\item In N-of-1 clinical trial simulation, Hendrickson et al.
  (2020) appears unique in using it.
\end{itemize}

\end{bullets}

\begin{rgt}
rgt to complete.

\end{rgt}

\begin{orig}
Architecture B (differential correlation via MVN) appears primarily in
latent variable and structural equation modeling contexts, including the
joint modeling framework discussed by
Rizopoulos\textsuperscript{\citeproc{ref-rizopoulos2012}{24}} (where
treatment-state-dependent covariance relates longitudinal biomarkers to
survival outcomes rather than to treatment response directly), Bayesian
joint modeling frameworks, and copula-based simulation studies. In the
clinical trial simulation literature specifically, the Hendrickson et
al. (2020) framework appears to be unique in using differential
correlation as the mechanism for biomarker-treatment interaction in an
N-of-1 trial context.

\end{orig}

\begin{bullets}

\textbf{The near-universal use of Architecture A means most
published power estimates are optimistic for Architecture B
biomarkers under carryover.}

\begin{itemize}\setlength{\itemsep}{2pt}\setlength{\parskip}{0pt}
\item Virtually all published power estimates for
  biomarker-moderated trials assume Architecture A.
\item These estimates may be optimistic for biomarkers that
  operate through Architecture B (covariance-structure
  mechanism).
\item The additional carryover-induced power loss under
  Architecture B is not accounted for.
\end{itemize}

\end{bullets}

\begin{rgt}
rgt to complete.

\end{rgt}

\begin{orig}
This predominance of Architecture A has an important implication: the
power estimates reported in the vast majority of biomarker-moderated
trial simulation studies may be optimistic for biomarkers that operate
through an Architecture B mechanism, because these studies do not
account for the additional power loss that carryover produces when the
interaction signal resides in the covariance structure rather than the
mean structure.

\end{orig}

\subsection{7.3 Crossover and N-of-1 trial methodology}

\begin{bullets}

\textbf{The crossover trial literature addresses carryover as a
nuisance parameter but ignores the architectural distinction;
the standard recommendation applies differently under each.}

\begin{itemize}\setlength{\itemsep}{2pt}\setlength{\parskip}{0pt}
\item Senn (2002), Jones and Kenward (2014), and Sturdevant and
  Lumley (2021) treat carryover as a nuisance parameter; none
  engages with DGP architecture.
\item Standard recommendation (include carryover in model or
  use adequate washout) applies to both architectures.
\item Under Architecture B, crossover designs with short
  washout periods may be less suitable than previously
  recognized.
\end{itemize}

\end{bullets}

\begin{rgt}
rgt to complete.

\end{rgt}

\begin{orig}
The crossover trial literature, represented most thoroughly by
Senn\textsuperscript{\citeproc{ref-senn2002}{25}} and Jones and
Kenward\textsuperscript{\citeproc{ref-jones2014crossover}{26}},
addresses carryover extensively but does not distinguish between
architectures for biomarker-treatment interaction modeling. More recent
methodological work by Sturdevant and
Lumley\textsuperscript{\citeproc{ref-sturdevant2021carryover}{27}} on
testing for carryover effects via mixed-effects models likewise treats
carryover as a nuisance parameter to be estimated and adjusted for,
without engaging with how carryover interacts with biomarker-treatment
effect signals. The standard recommendation, namely to always include
carryover in the analysis model or to use adequate washout, applies to
both architectures but has different power implications under each. The
finding that Architecture B is more sensitive to carryover suggests that
crossover designs with short washout periods may be less suitable for
detecting correlation-based biomarker interactions than previously
recognized.

\end{orig}

\subsection{7.4 Precision medicine trial design}

\begin{bullets}

\textbf{Enrichment and biomarker-stratified designs assume
Architecture A for power calculations; the architectural choice
matters differently across three design families.}

\begin{itemize}\setlength{\itemsep}{2pt}\setlength{\parskip}{0pt}
\item Simon (2010) and Freidlin and Korn (2014) frame power
  calculations in the Architecture A regression idiom.
\item Three design families warrant separate consideration:
  adaptive enrichment, treatment-switching, and basket/umbrella
  trials.
\item The architectural assumption is rarely made explicit, so the
  practical implications are easy to overlook.
\end{itemize}

\end{bullets}

\begin{rgt}
rgt to complete.

\end{rgt}

\begin{orig}
Enrichment and biomarker-stratified designs, as discussed by
Simon\textsuperscript{\citeproc{ref-simon2010}{5}} and Freidlin and
Korn\textsuperscript{\citeproc{ref-freidlin2014}{6}}, use Architecture A
exclusively for power calculations. Three families of design are
particularly worth considering in light of the architectural
distinction.

\end{orig}

\begin{bullets}

\textbf{Adaptive enrichment trials are vulnerable to
Architecture B power shortfalls precisely when the biomarker
mechanism is most uncertain.}

\begin{itemize}\setlength{\itemsep}{2pt}\setlength{\parskip}{0pt}
\item Interim power calculations driving adaptive decisions use
  Architecture A; if Architecture B holds and off-drug periods occur,
  realized power will be lower than predicted.
\item The mismatch biases adaptation rules toward premature
  futility stopping or overly narrow enrichment criteria.
\item The problem is worst when the biomarker mechanism is
  unknown, exactly when adaptive enrichment is most often
  proposed.
\end{itemize}

\end{bullets}

\begin{rgt}
rgt to complete.

\end{rgt}

\begin{orig}
\textbf{Adaptive enrichment trials} adjust the eligible population
mid-trial based on emerging evidence of a biomarker-treatment
interaction. The interim power calculations that drive these adaptive
decisions are universally framed in the Architecture A regression idiom.
If the true biomarker mechanism is closer to Architecture B and any
participants pass through washout intervals between dose-finding stages,
the realized power for the interaction test will be lower than the
adaptation rule expects, biasing decisions toward premature futility
stopping or toward overly narrow enrichment criteria. The discrepancy is
most acute when the biomarker mechanism is uncertain at the design
stage, which is precisely the situation in which adaptive enrichment is
most attractive.

\end{orig}

\begin{bullets}

\textbf{Treatment-switching designs suffer Architecture B power
loss through off-drug correlation erosion; parallel-group enrichment designs
do not, because there is no off-drug phase.}

\begin{itemize}\setlength{\itemsep}{2pt}\setlength{\parskip}{0pt}
\item Crossover, basket, platform, and N-of-1 hybrid designs all
  cycle participants through drug states; each off-drug phase
  erodes BM-BR correlation under Architecture B.
\item Parallel-group enrichment designs escape this loss because
  each participant stays on a single treatment throughout.
\item The architectural choice matters most in within-subject
  designs, where it is also most likely to be clinically
  important.
\end{itemize}

\end{bullets}

\begin{rgt}
rgt to complete.

\end{rgt}

\begin{orig}
\textbf{Treatment-switching designs} (most crossover, basket, and
platform trials, as well as N-of-1 hybrid designs) cycle individual
participants through different drug states. Under Architecture B, each
off-drug phase erodes the BM-BR correlation, so the biomarker-treatment
interaction signal in the second and subsequent treatment periods is
structurally weaker than in the first. The power loss documented in
Section 3 applies here directly. By contrast, parallel-group enrichment
designs, in which each participant remains on a single treatment for the
duration of the trial, escape this loss because there is no off-drug
phase during which the correlation can decay; for these designs,
Architectures A and B yield similar power estimates. The architectural
choice therefore matters most in exactly the designs where the biomarker
mechanism is most likely to be of clinical interest, namely those that
deliberately expose each participant to both treatment states in order
to estimate within-subject responsiveness.

\end{orig}

\begin{bullets}

\textbf{Basket and umbrella trials face an additional risk of
optimistic aggregated power estimates when Architecture B holds
in even a subset of arms.}

\begin{itemize}\setlength{\itemsep}{2pt}\setlength{\parskip}{0pt}
\item Renfro and Sargent (2017) discuss power calculations that
  average across biomarker-defined arms; the calculation assumes
  Architecture A throughout.
\item Mixed architectures across baskets produce optimistic
  trial-level power estimates that are difficult to detect from
  aggregate results.
\item Subgroup-specific power diagnostics under both
  architectures are warranted when the biological mechanism is
  heterogeneous.
\end{itemize}

\end{bullets}

\begin{rgt}
rgt to complete.

\end{rgt}

\begin{orig}
\textbf{Basket and umbrella trials}, as discussed by Renfro and
Sargent\textsuperscript{\citeproc{ref-renfro2017}{7}}, introduce an
additional layer because their power calculations average across
multiple biomarker-defined arms. If even a subset of the arms operate
through an Architecture B mechanism while the trial-wide power
calculation assumes Architecture A throughout, the aggregated power
estimates will be biased optimistically in a way that is difficult to
detect from trial-level results alone. Subgroup-specific power
diagnostics under both architectures should accompany the primary
planning analysis whenever the underlying biology is heterogeneous
across baskets.

\end{orig}

\begin{bullets}

\textbf{Across all three design families, the practical
recommendation is to anchor sample size planning on the
Architecture B estimate when the biomarker mechanism is
uncertain.}

\begin{itemize}\setlength{\itemsep}{2pt}\setlength{\parskip}{0pt}
\item Run power calculations under both architectures and treat
  Architecture B as the conservative anchor.
\item This recommendation parallels the advice in Section 5.3
  and is consistent with the PATH Statement (2020).
\item The PATH Statement's caution about data-driven
  interaction discovery reinforces the case for conservative
  planning under architectural uncertainty.
\end{itemize}

\end{bullets}

\begin{rgt}
rgt to complete.

\end{rgt}

\begin{orig}
In all three settings the practical recommendation parallels that of
Section 5.3: when the biological mechanism is uncertain, run the power
calculation under both architectures and treat the Architecture B
estimate as the more conservative anchor for sample size planning. This
recommendation is consistent with the PATH
Statement\textsuperscript{\citeproc{ref-kent2020path}{4}}, which itself
counsels caution about data-driven discovery of treatment-by-covariate
interactions and prefers risk modeling approaches when the application
criteria are met.

\end{orig}

\bigskip\noindent

\rule{\textwidth}{0.4pt}\bigskip

\section{8. Conclusions}

In conclusion, six points are to be emphasized.

\begin{enumerate}
\item The choice between direct mean moderation (Architecture A) and
  differential correlation (Architecture B) for generating
  biomarker-treatment interactions in clinical trial simulations
  has substantive consequences for power estimates under
  carryover.

\item Under Architecture A, carryover reduces power only
  modestly, by 1 to 3\% in relative terms across the three
  designs examined, because the proportional biomarker-drug
  relationship is preserved during the off-drug period. Under
  Architecture B the loss is strongly design-dependent: it
  remains modest in the crossover and hybrid designs (3 to
  5\%) but reaches roughly 31\% in the open-label OL+BDC
  design, because the differential-correlation signal erodes
  as the drug effect wanes.

\item The choice should be guided by the hypothesized biological
  mechanism: Architecture A for biomarkers that directly
  determine drug pharmacokinetics or pharmacodynamics;
  Architecture B for biomarkers that statistically predict
  treatment response through latent subtype membership.

\item For the prazosin-PTSD application with blood pressure as the
  predictive biomarker, Architecture B (MVN differential
  correlation) is the more appropriate model, and the resulting
  power sensitivity to carryover should inform trial design
  decisions regarding off-drug period duration.

\item The existing clinical trial simulation literature uses
  Architecture A almost exclusively. The Hendrickson et al.
  (2020) MVN approach is, to the best of our knowledge, unique
  in the N-of-1 trial context. This means that published power
  estimates for biomarker-moderated crossover and N-of-1
  designs may be systematically optimistic for biomarkers
  that operate through a correlation-based (Architecture B)
  mechanism, because they do not account for the carryover
  sensitivity documented here.

\item Simulation studies for biomarker-moderated treatment effects
  should explicitly report their DGP architecture and consider
  sensitivity analyses under the alternative when the biological
  mechanism is uncertain.
\end{enumerate}

\bigskip\noindent

\rule{\textwidth}{0.4pt}\bigskip

\section{Conflict of interest}

The authors declare no conflicts of interest.

\section{Funding}

This work received no specific funding.

\section{Data availability statement}

The data and code that support the findings of this study are openly
available in the pmsimstats-ng project repository. Per-replicate Monte
Carlo outputs and ADEMP-compliant cell-level summaries are versioned
under \texttt{analysis/data/}; simulation drivers are at
\texttt{analysis/scripts/}. Exact paths for this manuscript are listed
in the Reproducibility section below.

\bigskip\noindent

\rule{\textwidth}{0.4pt}\bigskip

\section{Reproducibility}

This manuscript was rendered with R\textasciitilde4.6.0 inside the
project's zzcollab Docker container (\texttt{rocker/tidyverse:4.6.0};
image hash recorded in \texttt{Dockerfile}; package set captured in
\texttt{renv.lock}). Principal packages used by the simulation pipeline
are \texttt{nlme} (continuous-time linear-mixed analysis with
\texttt{corCAR1} residual correlation), \texttt{parallel} with
\texttt{mclapply} (Section\textasciitilde3.1 and 3.3 production runs),
\texttt{data.table} (the \texttt{original-extended} simulation
collection), \texttt{purrr} (the \texttt{tidyverse} simulation
collection), and \texttt{rmarkdown} for the manuscript build.

\textbf{Sample-size convention.} The Section\textasciitilde3.1 driver
(\texttt{01-dgp-prototype.R}) generates \(N = 35\) subjects
independently per randomization path, then concatenates paths before
analysis; for two-path designs (CO, OL+BDC) the effective analysis
sample is therefore 70 subjects, and for the four-path Hybrid design it
is 140 subjects.

\textbf{Random-number generators.} The Section\textasciitilde3.1 driver
sets \texttt{set.seed(20260509)} and uses per-replicate seeds derived as
\texttt{100000 * cell\_index + rep\_index}.

\textbf{Data and code availability.} Sections 3.1 and 3.3: per-replicate
outputs at \texttt{analysis/data/quick-sim/01-dgp-replicates.rds}; ADEMP
cell-level summaries at \texttt{01-dgp-summary.txt}; driver at
\texttt{analysis/scripts/quick-sim/01-dgp-prototype.R}. The manuscript
source, bibliography, and CSL file are at
\texttt{analysis/report/01-dgp-mean-moderation-vs-mvn/}.

\section{References}\label{references}

\protect\phantomsection\label{refs}
\begin{CSLReferences}{0}{1}
\bibitem[\citeproctext]{ref-pmsimstats-paper10}
\CSLLeftMargin{1. }%
\CSLRightInline{{pmsimstats team}. Type {I} calibration of fixed-effect
interaction tests in linear mixed models: A staged diagnosis of
working-covariance misspecification and a cluster-robust remedy.
Published online 2026.}

\bibitem[\citeproctext]{ref-pmsimstats-paper08}
\CSLLeftMargin{2. }%
\CSLRightInline{{pmsimstats team}. Test-procedure choice for the
biomarker-treatment interaction in aggregated {N}-of-1 trials:
Dichotomization, classical {RM-ANOVA}, mixed-effects, and generalized
estimating equations. Published online 2026.}

\bibitem[\citeproctext]{ref-rothwell2005}
\CSLLeftMargin{3. }%
\CSLRightInline{Rothwell PM. Treating individuals 5. Subgroup analysis
in randomised controlled trials: Importance, indications, and
interpretation. \emph{The Lancet}. 2005;365(9454):176-186.}

\bibitem[\citeproctext]{ref-kent2020path}
\CSLLeftMargin{4. }%
\CSLRightInline{{Kent DM, Paulus JK, Klaveren D van, et al.} The
predictive approaches to treatment effect heterogeneity ({PATH})
statement. \emph{Annals of Internal Medicine}. 2020;172(1):35-45.}

\bibitem[\citeproctext]{ref-simon2010}
\CSLLeftMargin{5. }%
\CSLRightInline{Simon R. Clinical trial designs for evaluating the
medical utility of prognostic and predictive biomarkers in oncology.
\emph{Personalized Medicine}. 2010;7(1):33-47.}

\bibitem[\citeproctext]{ref-freidlin2014}
\CSLLeftMargin{6. }%
\CSLRightInline{Freidlin B, Korn EL. Biomarker enrichment strategies:
Matching trial design to biomarker credentials. \emph{Nature Reviews
Clinical Oncology}. 2014;11(2):81-90.}

\bibitem[\citeproctext]{ref-renfro2017}
\CSLLeftMargin{7. }%
\CSLRightInline{Renfro LA, Sargent DJ. Statistical controversies in
clinical research: Basket trials, umbrella trials, and other master
protocols. \emph{Annals of Oncology}. 2017;28(1):34-43.}

\bibitem[\citeproctext]{ref-haller2018}
\CSLLeftMargin{8. }%
\CSLRightInline{Haller B, Ulm K. A simulation study on estimating
biomarker-treatment interaction effects in randomized trials with
prognostic variables. \emph{Trials}. 2018;19(1):128.}

\bibitem[\citeproctext]{ref-zucker1997}
\CSLLeftMargin{9. }%
\CSLRightInline{Zucker DR, Schmid CH, McIntosh MW, D'Agostino RB, Selker
HP, Lau J. Combining single patient ({N}-of-1) trials to estimate
population treatment effects and to evaluate individual patient
responses to treatment. \emph{Journal of Clinical Epidemiology}.
1997;50(4):401-410.}

\bibitem[\citeproctext]{ref-araujo2016}
\CSLLeftMargin{10. }%
\CSLRightInline{Araujo A, Julious S, Senn S. Understanding variation in
sets of {N}-of-1 trials. \emph{PLoS ONE}. 2016;11(12):e0167167.}

\bibitem[\citeproctext]{ref-wang2019}
\CSLLeftMargin{11. }%
\CSLRightInline{Wang Y, Schork NJ. Power and design issues in
crossover-based {N}-of-1 clinical trials with fixed data collection
periods. \emph{Healthcare (Basel)}. 2019;7(3):84.}

\bibitem[\citeproctext]{ref-schork2022}
\CSLLeftMargin{12. }%
\CSLRightInline{Schork NJ. Accommodating serial correlation and
sequential design elements in personalized studies and aggregated
personalized studies. \emph{Harvard Data Science Review}. 2022;4(SI3).
doi:\href{https://doi.org/10.1162/99608f92.c4b04c64}{10.1162/99608f92.c4b04c64}}

\bibitem[\citeproctext]{ref-hendrickson2020}
\CSLLeftMargin{13. }%
\CSLRightInline{Hendrickson RC, Thomas RG, Schork NJ, Raskind MA.
Optimizing aggregated {N}-of-1 trial designs for predictive biomarker
validation: Statistical methods and practical considerations.
\emph{Frontiers in Digital Health}. 2020;2:13.
doi:\href{https://doi.org/10.3389/fdgth.2020.00013}{10.3389/fdgth.2020.00013}}

\bibitem[\citeproctext]{ref-pmsimstats-paper03}
\CSLLeftMargin{14. }%
\CSLRightInline{{pmsimstats team}. Latent-class and mixture-model
formulations for biomarker-treatment interaction in {N}-of-1 trials.
Published online 2026.}

\bibitem[\citeproctext]{ref-raskind2013}
\CSLLeftMargin{15. }%
\CSLRightInline{{Raskind MA, Peterson K, Williams T, et al.} A trial of
prazosin for combat trauma {PTSD} with nightmares in active-duty
soldiers returned from iraq and afghanistan. \emph{American Journal of
Psychiatry}. 2013;170(9):1003-1010.}

\bibitem[\citeproctext]{ref-raskind2018}
\CSLLeftMargin{16. }%
\CSLRightInline{{Raskind MA, Peskind ER, Chow B, et al.} Trial of
prazosin for post-traumatic stress disorder in military veterans.
\emph{New England Journal of Medicine}. 2018;378(6):507-517.}

\bibitem[\citeproctext]{ref-pmsimstats-paper11}
\CSLLeftMargin{17. }%
\CSLRightInline{{pmsimstats team}. A combined data-generating
architecture for biomarker-treatment interactions: Channel
super-additivity and mixing-weight sensitivity in aggregated {N}-of-1
trials. Published online 2026.}

\bibitem[\citeproctext]{ref-raison2013}
\CSLLeftMargin{18. }%
\CSLRightInline{{Raison CL, Rutherford RE, Woolwine BJ, et al.} A
randomized controlled trial of the tumor necrosis factor antagonist
infliximab for treatment-resistant depression: The role of baseline
inflammatory biomarkers. \emph{JAMA Psychiatry}. 2013;70(1):31-41.}

\bibitem[\citeproctext]{ref-senn2016mastering}
\CSLLeftMargin{19. }%
\CSLRightInline{Senn S. Mastering variation: Variance components and
personalised medicine. \emph{Statistics in Medicine}.
2016;35(7):966-977.
doi:\href{https://doi.org/10.1002/sim.6739}{10.1002/sim.6739}}

\bibitem[\citeproctext]{ref-raskind2003}
\CSLLeftMargin{20. }%
\CSLRightInline{{Raskind MA, Peskind ER, Kanter ED, et al.} Reduction of
nightmares and other {PTSD} symptoms in combat veterans by prazosin: A
placebo-controlled study. \emph{American Journal of Psychiatry}.
2003;160(2):371-373.}

\bibitem[\citeproctext]{ref-raskind2007}
\CSLLeftMargin{21. }%
\CSLRightInline{{Raskind MA, Peskind ER, Hoff DJ, et al.} A parallel
group placebo-controlled study of prazosin for trauma nightmares and
sleep disturbance in combat veterans with post-traumatic stress
disorder. \emph{Biological Psychiatry}. 2007;61(8):928-934.}

\bibitem[\citeproctext]{ref-morris2019simulation}
\CSLLeftMargin{22. }%
\CSLRightInline{Morris TP, White IR, Crowther MJ. Using simulation
studies to evaluate statistical methods. \emph{Statistics in Medicine}.
2019;38(11):2074-2102.
doi:\href{https://doi.org/10.1002/sim.8086}{10.1002/sim.8086}}

\bibitem[\citeproctext]{ref-duan2013}
\CSLLeftMargin{23. }%
\CSLRightInline{Duan N, Kravitz RL, Schmid CH. Single-patient ({N}-of-1)
trials: A pragmatic clinical decision methodology. \emph{Journal of
Clinical Epidemiology}. 2013;66(8):S21-S28.}

\bibitem[\citeproctext]{ref-rizopoulos2012}
\CSLLeftMargin{24. }%
\CSLRightInline{Rizopoulos D. \emph{Joint Models for Longitudinal and
Time-to-Event Data: With Applications in {R}}. CRC Press; 2012.}

\bibitem[\citeproctext]{ref-senn2002}
\CSLLeftMargin{25. }%
\CSLRightInline{Senn S. \emph{Cross-over Trials in Clinical Research}.
2nd ed. John Wiley \& Sons; 2002.}

\bibitem[\citeproctext]{ref-jones2014crossover}
\CSLLeftMargin{26. }%
\CSLRightInline{Jones B, Kenward MG. \emph{Design and Analysis of
Cross-over Trials}. 3rd ed. CRC Press; 2014.}

\bibitem[\citeproctext]{ref-sturdevant2021carryover}
\CSLLeftMargin{27. }%
\CSLRightInline{Sturdevant SG, Lumley T. Statistical methods for testing
carryover effects: A mixed effects model approach. \emph{Contemporary
Clinical Trials Communications}. 2021;22:100711.
doi:\href{https://doi.org/10.1016/j.conctc.2021.100711}{10.1016/j.conctc.2021.100711}}

\end{CSLReferences}

\vfill

\begin{center}\rule{0.5\linewidth}{0.5pt}\end{center}

\emph{Rendered on 2026-07-31 at 11:48 PDT.}\\
\emph{Source:
\texttt{\textasciitilde{}/prj/res/36-pmsimstats-ng/pmsimstats-ng/analysis/report/01-dgp-mean-moderation-vs-mvn/archive/report.Rmd}}

\end{document}
